% ================================================================
%  Rapport de Projet de Fin d'Année
%  Système de Recommandation Proactif ONCF -- Zero-Click Search
%  Auteur : Omar Chekroun
%  Etablissement : Université Mohammed V de Rabat
%                  Faculté des Sciences de Rabat
%  Année universitaire : 2025 -- 2026
%  Compilateur : pdfLaTeX (Overleaf)
% ================================================================
\documentclass[11pt, a4paper]{article}
\usepackage[utf8]{inputenc}
\usepackage[T1]{fontenc}
\usepackage[french]{babel}
\usepackage{geometry}
\usepackage{graphicx}
\usepackage{hyperref}
\usepackage{booktabs}
\usepackage{tabularx}
\usepackage{xcolor}
\usepackage{fancyhdr}
\usepackage{amsmath, amssymb}
\usepackage{listings}
\usepackage{enumitem}
\usepackage{multirow}
\usepackage{caption}
\usepackage{titlesec}
\usepackage{setspace}
\usepackage{array}

\geometry{margin=1in, top=0.9in, bottom=0.9in}
\onehalfspacing

\pagestyle{fancy}
\fancyhf{}
\rhead{\small Système de Recommandation Proactif ONCF}
\lhead{\small Rapport de PFA --- Omar Chekroun}
\cfoot{\thepage}

\titleformat{\section}
  {\Large\bfseries}{\thesection}{1em}{}
\titleformat{\subsection}
  {\large\bfseries}{\thesubsection}{1em}{}
\titleformat{\subsubsection}
  {\normalsize\bfseries\itshape}{\thesubsubsection}{1em}{}

\definecolor{codegray}{gray}{0.95}
\lstset{
  backgroundcolor=\color{codegray},
  basicstyle=\ttfamily\small,
  breaklines=true,
  frame=single,
  captionpos=b,
  tabsize=4,
  showstringspaces=false,
  keywordstyle=\bfseries,
  commentstyle=\itshape\color{gray},
  literate={é}{{\'{e}}}1{è}{{\`{e}}}1{ê}{{\^{e}}}1{à}{{\`{a}}}1
           {ù}{{\`{u}}}1{ô}{{\^{o}}}1{î}{{\^{i}}}1{ç}{{\c{c}}}1
           {É}{{\'{E}}}1{È}{{\`{E}}}1{À}{{\`{A}}}1,
}
\lstdefinestyle{py}{language=Python,
  morekeywords={True,False,None,self,dataclass,async,await}}
\lstdefinestyle{sh}{language=bash}
\lstdefinestyle{jn}{language={},morestring=[b]"}

\captionsetup{font=small,labelfont=bf}

% ================================================================
\begin{document}
% ================================================================

% ─────────────────────── PAGE DE GARDE ──────────────────────────
\begin{titlepage}
\thispagestyle{empty}
\begin{center}

\vspace*{0.3cm}
{\large \textbf{ROYAUME DU MAROC}}\\[0.15cm]
{\normalsize Université Mohammed V de Rabat}\\[0.1cm]
{\normalsize Faculté des Sciences de Rabat}\\[0.1cm]
{\normalsize Département Informatique}\\[0.6cm]

\noindent\rule{\textwidth}{1.2pt}
\vspace{0.05cm}\\
\noindent\rule{\textwidth}{0.4pt}

\vspace{1.5cm}
{\small \textsc{Rapport de Projet de Fin d'Année}}\\[0.4cm]
{\small Filière~: Sciences de l'Ingénieur Informatique}

\vspace{2.0cm}
{\Huge \textbf{Système de Recommandation}}\\[0.3cm]
{\Huge \textbf{Proactif ONCF}}\\[0.6cm]
{\Large \textit{«~Zero-Click Search~»}}\\[0.3cm]
{\large Prédiction du trajet ferroviaire suivant pour}\\[0.1cm]
{\large l'application mobile ONCF Voyages}

\vspace{2.0cm}

\noindent\rule{\textwidth}{0.4pt}
\vspace{0.05cm}\\
\noindent\rule{\textwidth}{1.2pt}

\vspace{0.8cm}
\begin{tabularx}{0.85\textwidth}{|l|X|}
\hline
\textbf{Réalisé par}        & Omar \textsc{Chekroun} \\
\hline
\textbf{Établissement}      & UM5 --- Faculté des Sciences de Rabat \\
\hline
\textbf{Filière}            & Sciences de l'Ingénieur Informatique \\
\hline
\textbf{Cadre du projet}    & Office National des Chemins de Fer du Maroc \\
                            & Direction Digitale \& Innovation \\
\hline
\textbf{Contact}            & omarchekroun39@gmail.com \\
\hline
\textbf{Année universitaire}& 2025 -- 2026 \\
\hline
\end{tabularx}

\vfill
{\large \textbf{Soutenu en Mai 2026}}\\[0.4cm]
\noindent\rule{\textwidth}{1.2pt}
\vspace{0.05cm}\\
\noindent\rule{\textwidth}{0.4pt}

\end{center}
\end{titlepage}

% ─────────────────────── DÉDICACES ──────────────────────────────
\newpage
\thispagestyle{empty}
\vspace*{2cm}
\begin{flushright}
\textit{\large À mes parents,}\\[0.2cm]
\textit{pour leur soutien indéfectible et leurs sacrifices.}\\[1cm]

\textit{\large À mes enseignants,}\\[0.2cm]
\textit{pour la rigueur et la passion qu'ils m'ont transmises.}\\[1cm]

\textit{\large À mes camarades de promotion,}\\[0.2cm]
\textit{pour les heures de travail partagées.}\\[1cm]

\textit{\large À tous ceux qui rendent}\\[0.2cm]
\textit{le voyage plus fluide pour les autres.}\\[2cm]

\hfill\textit{Omar Chekroun}
\end{flushright}

% ─────────────────────── REMERCIEMENTS ──────────────────────────
\newpage
\thispagestyle{fancy}
\section*{Remerciements}
\addcontentsline{toc}{section}{Remerciements}

Au terme de ce Projet de Fin d'Année, je tiens à exprimer ma profonde gratitude
à toutes les personnes qui ont contribué, de près ou de loin, à la réussite de
ce travail.

Je remercie en premier lieu le corps enseignant du \textbf{Département
Informatique de la Faculté des Sciences de Rabat}, et plus largement
\textbf{l'Université Mohammed V de Rabat}, pour la qualité de la formation
dispensée tout au long de mon cursus. Les enseignements en algorithmique,
en bases de données, en intelligence artificielle et en génie logiciel ont
constitué le socle indispensable à la conduite de ce projet.

Mes remerciements s'adressent également à \textbf{l'Office National des Chemins
de Fer (ONCF)}, et en particulier à sa \textit{Direction Digitale et
Innovation}, pour avoir mis à disposition un cas d'usage industriel d'une
grande richesse. La possibilité de travailler sur des données réelles de
réservation, à l'échelle nationale, a transformé ce projet académique en une
expérience d'ingénierie authentique, soumise aux mêmes contraintes que celles
rencontrées en entreprise~: volumétrie, qualité hétérogène des données, exigences
de conformité réglementaire (Loi 09-08 / CNDP) et impératifs de performance.

Je remercie également la communauté open source~: les auteurs et mainteneurs des
bibliothèques \texttt{XGBoost}, \texttt{scikit-learn}, \texttt{pandas},
\texttt{FastAPI} et \texttt{Pydantic}, sans lesquels un tel pipeline n'aurait
pu être construit en quelques semaines par une seule personne.

Enfin, ma reconnaissance va à ma famille, dont le soutien moral et matériel
est resté constant durant toute la période de réalisation de ce projet.

\vspace{1cm}
\hfill\textit{Omar Chekroun}\\
\hfill\textit{Mai 2026}

% ─────────────────────── RÉSUMÉ ─────────────────────────────────
\newpage
\section*{Résumé}
\addcontentsline{toc}{section}{Résumé}

Ce rapport présente la conception, le développement et l'évaluation d'un
\textbf{système de recommandation proactif} de trajets ferroviaires destiné à
l'application mobile ONCF Voyages. Le projet adresse un problème métier
précis~: la \textit{friction de recherche} subie quotidiennement par les
voyageurs réguliers, contraints de saisir manuellement leurs trajets
récurrents. La solution proposée, désignée sous le terme
\textbf{«~Zero-Click Search~»}, prédit le trajet suivant le plus probable
dès l'ouverture de l'application et le présente à l'utilisateur sans aucune
interaction.

Le système est construit autour d'une \textbf{architecture deux étapes}
inspirée des pratiques d'Uber Eats et d'Airbnb~: une phase heuristique de
\textit{Candidate Generation} sélectionne dix liaisons plausibles à partir
de l'historique de l'utilisateur, et une phase de \textit{Ranking} par un
classifieur \textbf{XGBoost multiclasse} (1\,011 classes) ordonne ces candidats
selon leur probabilité d'être le prochain trajet. Cette séparation garantit
à la fois la scalabilité (filtrage rapide d'un large espace de classes) et la
qualité de prédiction (modèle entraîné spécifiquement pour départager les
candidats pertinents).

Le pipeline complet comprend~: (i)~le \textbf{nettoyage} de 946\,155 lignes
brutes de réservation, ramenées à 491\,680 lignes exploitables après
application de règles métier strictes (suppression des annulations propagées,
filtrage cold start, déduplication)~; (ii)~l'\textbf{ingénierie de 26
variables} comportementales et temporelles, incluant un encodage cyclique des
horodatages~; (iii)~l'\textbf{entraînement supervisé} sur un split temporel
80/20 (393\,344 lignes d'entraînement, 98\,261 lignes de test)~; (iv)~le
\textbf{déploiement} via une API REST FastAPI à faible latence.

Les résultats offline atteints dépassent significativement les seuils
initialement fixés~:
\textbf{Hit~Rate@1 = 76{,}3~\%}, \textbf{Hit~Rate@3 = 90{,}6~\%},
\textbf{MRR@3 = 82{,}8~\%}. Le modèle est capable de prédire correctement le
trajet suivant dès la première suggestion dans plus de trois quarts des cas,
sur un espace de plus de mille classes possibles.

Le rapport décrit en détail chaque étape de la conception, les choix
techniques et leur justification, les difficultés rencontrées (saturation de
VRAM GPU, fuites de variable d'identité, gestion des classes inédites au test)
ainsi que leurs résolutions, et propose des perspectives d'amélioration~:
features de popularité globale, réentraînement automatisé, A/B testing en
production.

\bigskip
\noindent\textbf{Mots-clés~:} système de recommandation, prédiction proactive,
Zero-Click Search, XGBoost multiclasse, Learning to Rank, Candidate
Generation, FastAPI, pipeline scikit-learn, Loi 09-08 / CNDP, ONCF.

% ─────────────────────── ABSTRACT ───────────────────────────────
\newpage
\section*{Abstract}
\addcontentsline{toc}{section}{Abstract}

This report presents the design, development, and evaluation of a
\textbf{proactive railway trip recommender system} for the ONCF Voyages
mobile application. The project addresses a concrete business problem~: the
\textit{search friction} faced daily by regular passengers, who must manually
re-enter their recurring trips. The proposed solution, called
\textbf{«~Zero-Click Search~»}, predicts the most likely next trip as soon as
the application opens, and presents it to the user with no interaction
required.

The system follows a \textbf{two-stage architecture} inspired by industry
references such as Uber Eats and Airbnb~: a heuristic
\textit{Candidate Generation} phase selects ten plausible routes from the
user's history, and a \textit{Ranking} phase based on a \textbf{multiclass
XGBoost classifier} (1\,011 classes) orders these candidates by their
probability of being the next trip. This separation ensures both scalability
(fast filtering across a large class space) and prediction quality (a model
trained specifically to discriminate between relevant candidates).

The complete pipeline includes (i)~the \textbf{cleaning} of 946\,155 raw
booking rows down to 491\,680 usable rows after applying strict business rules
(propagated cancellation removal, cold-start filtering, deduplication)~;
(ii)~the \textbf{engineering of 26 features}, both behavioral and temporal,
with cyclic encoding of timestamps~; (iii)~\textbf{supervised training} on an
80/20 chronological split~; (iv)~\textbf{deployment} as a low-latency FastAPI
REST service.

Achieved offline metrics significantly exceed the initially fixed thresholds~:
\textbf{Hit~Rate@1 = 76.3\,\%}, \textbf{Hit~Rate@3 = 90.6\,\%},
\textbf{MRR@3 = 82.8\,\%}. The model correctly predicts the next trip as the
first suggestion in nearly three out of four cases, over a space of more than
one thousand possible classes.

\bigskip
\noindent\textbf{Keywords~:} recommender system, proactive prediction,
Zero-Click Search, multiclass XGBoost, Learning to Rank, Candidate Generation,
FastAPI, scikit-learn pipeline, Moroccan Data Protection Law (09-08), ONCF.

% ─────────────────────── TABLE DES MATIÈRES ─────────────────────
\newpage
\tableofcontents

% ─────────────────────── LISTE DES FIGURES ──────────────────────
\newpage
\section*{Liste des Tableaux}
\addcontentsline{toc}{section}{Liste des Tableaux}
\listoftables

% ─────────────────────── LISTE DES ABRÉVIATIONS ─────────────────
\newpage
\section*{Liste des Abréviations}
\addcontentsline{toc}{section}{Liste des Abréviations}

\begin{tabularx}{\textwidth}{|l|X|}
\hline
\textbf{Acronyme} & \textbf{Signification} \\
\hline
ONCF        & Office National des Chemins de Fer (du Maroc) \\
\hline
CNDP        & Commission Nationale de contrôle de la protection des Données à caractère Personnel \\
\hline
PFA         & Projet de Fin d'Année \\
\hline
UM5         & Université Mohammed V de Rabat \\
\hline
ML          & Machine Learning (apprentissage automatique) \\
\hline
GBDT        & Gradient Boosted Decision Trees \\
\hline
XGBoost     & eXtreme Gradient Boosting (bibliothèque de GBDT) \\
\hline
API         & Application Programming Interface \\
\hline
REST        & REpresentational State Transfer \\
\hline
HTTP        & HyperText Transfer Protocol \\
\hline
JSON        & JavaScript Object Notation \\
\hline
HR@k        & Hit Rate at $k$ (taux de présence dans le top-$k$) \\
\hline
MRR         & Mean Reciprocal Rank (rang réciproque moyen) \\
\hline
O/D         & Origine / Destination \\
\hline
OHE         & One-Hot Encoding \\
\hline
OE          & Ordinal Encoding \\
\hline
OOM         & Out Of Memory (épuisement mémoire) \\
\hline
GPU         & Graphics Processing Unit \\
\hline
CPU         & Central Processing Unit \\
\hline
VRAM        & Video Random Access Memory \\
\hline
CUDA        & Compute Unified Device Architecture (NVIDIA) \\
\hline
CSV         & Comma-Separated Values \\
\hline
CTR         & Click-Through Rate (taux de clic) \\
\hline
\end{tabularx}

\newpage

% ================================================================
\section{Introduction Générale}
% ================================================================

\subsection{Contexte du Projet}

À l'ère du numérique, les utilisateurs d'applications mobiles attendent
des services qui anticipent leurs besoins plutôt que de simplement y
répondre. Cette tendance, observable chez les acteurs du e-commerce,
de la mobilité urbaine, de la livraison à la demande et de la diffusion
de contenu, est désormais un standard implicite~: l'utilisateur ouvre
l'application et se voit immédiatement proposer ce qu'il est
\textit{statistiquement le plus susceptible de chercher}. Cette
proactivité n'est pas un gadget~; elle réduit la \textit{friction de
recherche}, accélère la conversion et augmente la satisfaction perçue.

L'Office National des Chemins de Fer du Maroc (ONCF), opérateur public
en charge du réseau ferroviaire national, exploite une application
mobile dénommée \textit{ONCF Voyages}. Cette application permet aux
usagers de consulter les horaires, de réserver des billets et de gérer
leurs voyages. Or, à ce jour, l'expérience utilisateur reste largement
\textbf{réactive}~: à chaque ouverture, le voyageur doit ressaisir
manuellement les gares de départ et d'arrivée, sélectionner une date,
puis lancer une recherche. Pour les voyageurs occasionnels, ce processus
est acceptable. Pour les \textbf{voyageurs réguliers} --- qui empruntent
quasi quotidiennement les mêmes liaisons (typiquement domicile/travail
ou domicile/études) --- cette répétition manuelle est une perte de
temps et une source de frustration.

C'est précisément à cette catégorie d'utilisateurs que ce projet
s'adresse, à travers la conception d'un \textbf{système de
recommandation proactif} appelé \textbf{«~Zero-Click Search~»}.

\subsection{Problématique}

Le problème posé peut s'énoncer de manière concise~:
\begin{quote}
\textit{Étant donné l'historique complet des réservations d'un client,
ainsi que le contexte temporel d'ouverture de l'application, prédire la
prochaine liaison ONCF que ce client réservera, et lui présenter cette
prédiction (ou un top-3) sans qu'il ait à effectuer la moindre
recherche.}
\end{quote}

Cette problématique soulève plusieurs sous-questions techniques et métier
non triviales~:
\begin{enumerate}
\item \textbf{Représentation de l'utilisateur~:} comment résumer
    l'historique d'un voyageur (parfois plusieurs centaines de
    réservations) en une représentation exploitable par un modèle~?
\item \textbf{Espace de prédiction~:} le réseau ONCF compte plus de
    1\,000 liaisons distinctes après nettoyage. Comment formuler la
    prédiction parmi un espace aussi large~?
\item \textbf{Démarrage à froid (cold start)~:} que faire pour les
    nouveaux clients sans historique exploitable~? Risque-t-on de
    dégrader leur expérience en leur proposant des trajets non
    pertinents~?
\item \textbf{Conformité réglementaire~:} le Maroc applique la
    \textit{Loi 09-08} relative à la protection des données à caractère
    personnel, sous le contrôle de la CNDP. Quelle est la frontière
    entre identifiant de lookup et variable de modèle~?
\item \textbf{Latence et coût d'inférence~:} l'application mobile attend
    une réponse immédiate à chaque ouverture. Comment garantir une
    latence inférieure à quelques dizaines de millisecondes~?
\end{enumerate}

L'ensemble de ces questions a structuré les choix techniques de ce
projet, dont le présent rapport rend compte.

\subsection{Objectifs}

Les objectifs initiaux du projet, fixés en début de période, étaient
les suivants~:
\begin{itemize}
\item Concevoir une architecture de recommandation \textbf{compatible
    production}~: scalable, interprétable, à faible latence.
\item Atteindre des métriques offline minimales~: \textbf{Hit~Rate@1
    $> 30\,\%$}, \textbf{Hit~Rate@3 $> 50\,\%$}, \textbf{MRR@3 $>
    35\,\%$} sur un split temporel.
\item Respecter \textbf{strictement} la Loi 09-08~: aucune utilisation
    du \texttt{CodeClient} comme variable de modèle, aucune apparition
    dans les logs.
\item Livrer une \textbf{API REST opérationnelle} prête à être consommée
    par l'application mobile.
\item Documenter le projet de manière à permettre sa
    \textbf{reproduction complète} par un tiers.
\end{itemize}

À la date de rédaction de ce rapport, l'ensemble de ces objectifs est
atteint --- et les métriques obtenues dépassent largement les seuils
fixés (HR@1 = 76{,}3\,\% au lieu des 30\,\% attendus).

\subsection{Méthodologie de Travail}

Le projet a suivi une méthodologie structurée en cinq phases
chronologiques~:
\begin{enumerate}
\item \textbf{Phase d'étude (semaine 1--2)~:} analyse du besoin
    métier, étude de l'état de l'art, rédaction du benchmark, premiers
    choix d'architecture.
\item \textbf{Phase d'exploration des données (semaine 2--3)~:}
    réception et profilage du dataset brut, identification des problèmes
    de qualité, élaboration des règles métier de nettoyage.
\item \textbf{Phase de modélisation (semaine 3--6)~:} ingénierie
    des variables, sélection de l'algorithme, prototypage, premières
    expérimentations, ajustement des hyperparamètres.
\item \textbf{Phase de déploiement (semaine 6--7)~:} encapsulation
    du modèle dans un composant \texttt{Recommender}, exposition via une
    API REST FastAPI, écriture des tests unitaires.
\item \textbf{Phase de consolidation (semaine 7--8)~:} correction
    des bugs détectés, écriture du rapport, préparation de la
    soutenance.
\end{enumerate}

À chaque étape, des artefacts ont été produits et versionnés~: scripts
exécutables, rapports JSON, suite de tests automatisée, documentation
technique (\texttt{CLAUDE.md}, fichiers \texttt{README}).

\subsection{Plan du Rapport}

Le rapport s'articule en treize sections principales. La
\textbf{section~2} présente l'organisme d'accueil, l'ONCF, ainsi que sa
direction digitale. La \textbf{section~3} détaille l'étude de l'existant
et le benchmark des solutions industrielles. La \textbf{section~4}
formalise le cahier des charges. Les \textbf{sections~5 et~6} décrivent
les données et l'ingénierie des variables. La \textbf{section~7} présente
la conception du système et son architecture deux étapes. Les
\textbf{sections~8 et~9} approfondissent la modélisation XGBoost et les
difficultés rencontrées lors de l'entraînement. La \textbf{section~10}
décrit la méthodologie d'évaluation et les résultats obtenus. Les
\textbf{sections~11 et~12} traitent du déploiement (composant
\texttt{Recommender}, API REST) et de la qualité du code (tests
unitaires). La \textbf{section~13} dresse le bilan et trace les
perspectives. Enfin, plusieurs \textbf{annexes} regroupent les artefacts
techniques (rapports JSON, code source clé, commandes de reproduction).

\newpage
% ================================================================
\section{Présentation de l'Organisme d'Accueil}
% ================================================================

\subsection{L'ONCF~: Acteur Public Stratégique}

L'\textbf{Office National des Chemins de Fer (ONCF)} est l'établissement
public marocain en charge de l'exploitation, de la maintenance et du
développement du réseau ferroviaire national. Créé en 1963 par fusion de
trois compagnies privées historiques, l'ONCF est placé sous la tutelle du
Ministère de l'Équipement, du Transport et de la Logistique. Il assure
trois activités principales~: le \textit{transport de voyageurs}, le
\textit{transport de marchandises} et le \textit{développement
d'infrastructures ferroviaires}.

Le réseau ferroviaire marocain s'étend sur près de 2\,300 kilomètres,
reliant les principales villes du Royaume entre Tanger au nord et
Marrakech au sud, avec des liaisons transversales vers Oujda à l'est. La
mise en service en 2018 de la ligne à grande vitesse \textbf{Al~Boraq}
(Tanger--Casablanca, premier TGV du continent africain) a considérablement
accru la fréquentation et la complexité de l'offre commerciale.

\subsection{Volumétrie et Enjeux Numériques}

L'ONCF transporte annuellement plusieurs dizaines de millions de
voyageurs, ce qui en fait l'un des plus gros producteurs nationaux de
données de mobilité. La numérisation de la billetterie via l'application
mobile \textbf{ONCF~Voyages}, lancée puis enrichie au fil des années,
génère un \textbf{flux continu de données de réservation} structurées~:
identifiants client pseudonymisés, gares de départ et d'arrivée, dates
et heures de départ, classes, prix, canaux d'achat, etc.

Ces données constituent un actif stratégique de premier plan, à la fois
pour la planification de l'offre (programmation des trains, dimensionnement
des rames, gestion tarifaire) et pour l'amélioration continue de
l'expérience utilisateur. L'exploitation par des techniques d'apprentissage
automatique en est un axe naturel~: la régularité des comportements de
mobilité ferroviaire (trajets domicile/travail) se prête particulièrement
bien aux modèles prédictifs.

\subsection{Direction Digitale et Innovation}

La \textbf{Direction Digitale et Innovation} de l'ONCF est le service
porteur de la transformation numérique de l'établissement. Elle pilote
notamment~:
\begin{itemize}
\item l'évolution fonctionnelle de l'application mobile \textbf{ONCF
    Voyages} et du portail web de réservation~;
\item la mise en place de \textbf{services data} internes
    (entrepôts de données, pipelines analytiques)~;
\item l'expérimentation de \textbf{technologies d'intelligence
    artificielle} appliquées aux problématiques métier de l'opérateur~;
\item la \textbf{conformité réglementaire} en matière de protection des
    données personnelles (Loi 09-08).
\end{itemize}

C'est dans le cadre de la dernière mission citée --- l'expérimentation
de l'IA pour améliorer l'expérience utilisateur --- que s'inscrit le
présent projet de recommandation proactive.

\subsection{Cadre Réglementaire~: Loi 09-08 et CNDP}

La \textbf{Loi 09-08} (promulguée en 2009) encadre au Maroc le
traitement des données à caractère personnel. Elle est appliquée par la
\textbf{Commission Nationale de contrôle de la protection des Données à
caractère Personnel (CNDP)}. Toute solution traitant des données
nominatives ou indirectement identifiantes doit s'y conformer.

Pour le présent projet, deux principes directeurs sont retenus~:
\begin{enumerate}
\item Le \textbf{\texttt{CodeClient}} (identifiant interne ONCF) est
    utilisé \textit{uniquement comme clé de lookup} pour récupérer
    l'historique d'un utilisateur. Il \textit{n'est jamais utilisé comme
    variable d'apprentissage}~: si le modèle pouvait s'appuyer sur
    l'identifiant numérique, il développerait des prédictions par
    sur-apprentissage individuel impossible à généraliser et difficile à
    auditer.
\item Le \texttt{CodeClient} \textbf{n'apparaît dans aucun log de l'API}
    en clair, ni en réponse, ni en requête tracée. Les requêtes sont
    journalisées sans le champ \texttt{code\_client}.
\end{enumerate}

Ces deux règles ont été appliquées \textit{by design} dans le code, et
font l'objet de tests unitaires dédiés.

\newpage
% ================================================================
\section{Étude de l'Existant et Benchmark}
% ================================================================

\subsection{Pourquoi un Benchmark~?}

Avant d'engager un développement coûteux, il est méthodologiquement sain
d'examiner comment des problèmes similaires sont traités par les
acteurs leaders du marché. Cette étape, appelée \textit{benchmark},
poursuit trois objectifs~:
\begin{enumerate}
\item \textbf{Éviter la réinvention~:} si une approche éprouvée existe,
    autant s'en inspirer.
\item \textbf{Calibrer les attentes~:} comprendre les ordres de grandeur
    des performances atteignables et les contraintes opérationnelles
    typiques.
\item \textbf{Identifier les pièges~:} les retours d'expérience publics
    (articles de blog d'ingénierie, papiers de conférence, présentations
    en conférence) signalent souvent des écueils non triviaux.
\end{enumerate}

Quatre références ont été examinées en début de projet, choisies pour
leur diversité d'approche tout en partageant la nature
\textit{prédiction d'action utilisateur} avec notre cas d'usage.

\subsection{Référence 1~: Uber Eats / Uber Rides}

Uber publie régulièrement, sur son blog d'ingénierie, des descriptions
détaillées de ses systèmes de recommandation. L'architecture canonique
qui en ressort est le schéma \textbf{Candidate Generation +~Ranking}~:
\begin{itemize}
\item \textbf{Candidate Generation~:} pour chaque utilisateur, un
    ensemble de quelques dizaines (au plus quelques centaines) de
    candidats est extrait rapidement à partir de signaux simples~:
    historique récent, popularité locale, restaurants déjà commandés,
    courses fréquentes.
\item \textbf{Ranking~:} les candidats ainsi pré-filtrés sont ordonnés
    par un modèle XGBoost (ou plus récemment un réseau de neurones
    léger) entraîné sur des features comportementales et contextuelles
    plus riches.
\end{itemize}

\textbf{Leçon retenue~:} la séparation des deux phases est un design
pattern essentiel. Tenter de scorer un large espace de classes en une
passe directe est à la fois lent et statistiquement bruité. Le
pré-filtrage heuristique réduit l'espace tout en garantissant que les
candidats vraiment plausibles y figurent.

\subsection{Référence 2~: Trainline / Deutsche Bahn}

Trainline (agrégateur ferroviaire européen) et Deutsche Bahn
(opérateur historique allemand) déploient des recommandations basées
sur du \textbf{filtrage collaboratif} et des \textbf{règles métier}.
Le filtrage collaboratif s'appuie sur le principe~: \textit{les
utilisateurs qui ont réservé X ont aussi réservé Y}.

\textbf{Leçon retenue~:} le filtrage collaboratif est puissant pour les
réseaux denses où les comportements se chevauchent largement
(grandes villes européennes, dense maillage de gares). Il l'est
\textbf{moins} dans un contexte comme celui de l'ONCF, où les
habitudes ferroviaires sont très individuelles (un voyageur de Casablanca
fait rarement les mêmes trajets qu'un voyageur de Fès), et où le
comportement collectif risque de noyer le signal individuel. Ce
constat a écarté le filtrage collaboratif comme approche principale.

\subsection{Référence 3~: Airbnb}

Airbnb a publié plusieurs articles sur sa pile de recommandation, dont
le très remarqué \textit{Real-time Personalization using Embeddings for
Search Ranking at Airbnb} (KDD~2018). L'approche combine~:
\begin{itemize}
\item un modèle de \textbf{GBDT (Gradient Boosted Decision Trees)}
    pour le ranking final~;
\item des \textbf{embeddings appris} de destinations, capturant la
    similarité sémantique entre lieux (proche géographiquement, profil
    de voyageur similaire, saisonnalité comparable, etc.).
\end{itemize}

\textbf{Leçon retenue~:} les embeddings sont une voie d'enrichissement
puissante \textit{lorsque} la donnée est suffisamment dense. Le réseau
ferroviaire marocain présente certaines liaisons à très faible
fréquence (quelques voyages par mois), pour lesquelles l'apprentissage
d'embeddings serait peu fiable. Cette approche est mise de côté pour
une version~2 du projet, lorsque davantage de données auront été
collectées et qu'un graphe géographique enrichi sera disponible.

\subsection{Référence 4~: Google Maps}

Google Maps utilise des modèles probabilistes (chaînes de Markov
notamment) sur les séquences de déplacement pour anticiper les trajets
récurrents (\textit{driving home} le soir, etc.). Le signal exploité
est principalement la \textbf{fréquence} et la \textbf{récence}.

\textbf{Leçon retenue~:} ces deux signaux --- fréquence et récence
d'un trajet dans l'historique --- doivent être explicitement encodés
dans nos features. Ils constituent les deux variables ordinales les
plus discriminantes dans la phase de Candidate Generation.

\subsection{Synthèse Comparative}

\begin{table}[h]
\centering
\small
\begin{tabularx}{\textwidth}{|l|X|c|}
\hline
\textbf{Référence} & \textbf{Apport principal} & \textbf{Retenu~?} \\
\hline
Uber Eats          & Architecture deux étapes Candidate~+~Ranking & Oui (architecture) \\
\hline
Trainline / DB     & Filtrage collaboratif sur réseau dense       & Non (réseau peu dense) \\
\hline
Airbnb             & GBDT + embeddings de destinations            & Partiellement (GBDT oui, embeddings v2) \\
\hline
Google Maps        & Importance de fréquence + récence            & Oui (features) \\
\hline
\end{tabularx}
\caption{Synthèse du benchmark et choix de conception associés.}
\end{table}

\subsection{Décisions Architecturales Issues du Benchmark}

Trois décisions structurantes découlent directement de ce benchmark~:
\begin{enumerate}
\item \textbf{Architecture deux étapes (Uber).} On sépare la sélection
    des candidats (heuristique, déterministe, rapide) et leur classement
    (modèle XGBoost, probabiliste, plus riche).
\item \textbf{XGBoost multiclasse (Airbnb).} On retient un GBDT pour
    sa robustesse, sa rapidité, son interprétabilité (importance de
    variables exploitable) et sa compatibilité immédiate avec
    \texttt{scikit-learn}.
\item \textbf{Features fréquence/récence en pivot (Google Maps).} On
    encode explicitement, dans la table de features, l'index du voyage
    dans l'historique (\texttt{user\_trip\_index}), le délai depuis le
    voyage précédent (\texttt{days\_since\_prev}) et la liaison
    précédente (\texttt{prev\_liaison}).
\end{enumerate}

\newpage
% ================================================================
\section{Cahier des Charges et Spécifications}
% ================================================================

\subsection{Spécifications Fonctionnelles}

Le système doit répondre aux exigences fonctionnelles suivantes,
formulées en début de projet et révisées à mi-parcours~:

\begin{table}[h]
\centering
\small
\begin{tabularx}{\textwidth}{|l|l|X|}
\hline
\textbf{Code} & \textbf{Priorité} & \textbf{Description} \\
\hline
SF-01 & Haute & Fournir une recommandation de top-1 ou top-3 pour tout
                client donné, en moins de 100~ms par requête. \\
\hline
SF-02 & Haute & Si le client a moins de 3~réservations validées, retourner
                une réponse explicite \texttt{cold\_start} (sans deviner). \\
\hline
SF-03 & Haute & Ne jamais utiliser le \texttt{CodeClient} comme feature
                d'apprentissage. \\
\hline
SF-04 & Haute & Exposer une route HTTP \texttt{POST /recommend} qui accepte
                un \texttt{code\_client} et un entier $k \in [1, 3]$. \\
\hline
SF-05 & Moyenne & Exposer une route \texttt{GET /health} pour le monitoring
                  d'orchestrateur. \\
\hline
SF-06 & Moyenne & Documenter automatiquement l'API (Swagger / OpenAPI). \\
\hline
SF-07 & Basse  & Permettre la reconstruction complète du modèle depuis les
                 données brutes en exécutant trois scripts numérotés. \\
\hline
\end{tabularx}
\caption{Spécifications fonctionnelles du système.}
\end{table}

\subsection{Spécifications Non Fonctionnelles}

\begin{table}[h]
\centering
\small
\begin{tabularx}{\textwidth}{|l|X|}
\hline
\textbf{Code} & \textbf{Description} \\
\hline
SNF-01 & \textbf{Performance offline~:} HR@1 $> 0{,}30$, HR@3 $> 0{,}50$,
         MRR@3 $> 0{,}35$ sur split temporel 80/20. \\
\hline
SNF-02 & \textbf{Latence~:} temps de réponse de l'API $< 100$~ms en
         conditions normales (modèle déjà chargé en mémoire). \\
\hline
SNF-03 & \textbf{Empreinte mémoire~:} chargement du modèle $< 1$~Go RAM
         résidente. \\
\hline
SNF-04 & \textbf{Reproductibilité~:} déterminisme strict des étapes 1 à 3
         pour un même jeu de données et une même version des dépendances. \\
\hline
SNF-05 & \textbf{Conformité réglementaire~:} aucune fuite de
         \texttt{CodeClient} dans les artefacts ni dans les logs. \\
\hline
SNF-06 & \textbf{Maintenabilité~:} séparation stricte logique métier
         (\texttt{rec\_oncf}) / couche HTTP (\texttt{apps/api}). \\
\hline
SNF-07 & \textbf{Testabilité~:} couverture par tests unitaires sur les
         composants clés (au moins~25 tests, 100\,\% verts). \\
\hline
\end{tabularx}
\caption{Spécifications non fonctionnelles du système.}
\end{table}

\subsection{Périmètre Hors Scope}

Pour cadrer la durée du projet, plusieurs éléments sont \textbf{hors
scope} de cette première itération et explicitement renvoyés à des
versions ultérieures~:
\begin{itemize}
\item Intégration des données de \textbf{disponibilité temps réel} des
    trains (planning ONCF backend).
\item Recommandation \textbf{multi-trajets} (suggestion d'aller-retour
    cohérent, propositions de combinaisons).
\item \textbf{Personnalisation contextuelle géolocalisée} (utilisation
    du GPS pour choisir la gare la plus proche).
\item \textbf{Interface utilisateur mobile}~: ce projet livre l'API et
    le modèle, l'intégration en Android/iOS est laissée à l'équipe
    front.
\item \textbf{Réentraînement automatisé} (pipeline CI/CD de
    \textit{retraining} hebdomadaire avec promotion conditionnelle).
\end{itemize}

\subsection{Environnement Technique}

Le projet est développé et exécuté dans l'environnement suivant~:
\begin{itemize}
\item \textbf{Système d'exploitation~:} Windows 11 Pro (interpréteur
    PowerShell par défaut).
\item \textbf{Langage~:} Python 3.12.10.
\item \textbf{Bibliothèques principales~:} \texttt{pandas~3.x},
    \texttt{xgboost~3.2.0}, \texttt{scikit-learn},
    \texttt{joblib}, \texttt{fastapi~0.115}, \texttt{pydantic~v2},
    \texttt{uvicorn}.
\item \textbf{Matériel~:} ordinateur portable doté d'un GPU NVIDIA
    RTX~3050 (4~Go VRAM). À noter que le GPU n'a finalement pas été
    utilisé pour l'entraînement (cf. section~9.1).
\item \textbf{Environnement isolé~:} \texttt{venv} dans
    \texttt{.venv\textbackslash} à la racine du dépôt.
\end{itemize}

\newpage
% ================================================================
\section{Données~: Acquisition, Exploration, Nettoyage}
% ================================================================

\subsection{Sources de Données}

Deux fichiers \texttt{CSV} ont été fournis comme matière première~:

\begin{itemize}
\item \texttt{oncf\_data.csv}~: \textbf{table principale des
    réservations}. Une ligne par segment de voyage réservé. Volumétrie
    brute~: 946\,155 lignes, 18 colonnes.
\item \texttt{Liaison.csv}~: \textbf{table de référence des liaisons}.
    Une ligne par identifiant \texttt{LiaisonId}, avec les libellés
    associés des gares de départ et d'arrivée.
\end{itemize}

\subsection{Schéma Brut et Sémantique des Colonnes}

L'analyse du schéma brut a permis d'identifier la signification de
chaque colonne. Les colonnes essentielles sont décrites ci-dessous~:

\begin{table}[h]
\centering
\small
\begin{tabularx}{\textwidth}{|l|l|X|}
\hline
\textbf{Colonne} & \textbf{Type brut} & \textbf{Sémantique} \\
\hline
\texttt{CodeClient}                  & Numérique &
  Identifiant client interne ONCF (clé de lookup uniquement). \\
\hline
\texttt{AchteurId}                   & Numérique &
  Identifiant de l'acheteur (peut différer de CodeClient pour les achats
  pour autrui, professionnels, parents pour enfants). \\
\hline
\texttt{LiaisonId}                   & Mixte     &
  Identifiant de la liaison O/D. \textbf{Variable cible} du modèle. \\
\hline
\texttt{DateHeureDepartVoyageSegment}& Texte     &
  Horodatage du départ du segment de voyage. \\
\hline
\texttt{NbrVoySegment}               & Entier    &
  Nombre de voyageurs sur ce segment. Une valeur $\leq 0$ trahit une
  annulation. \\
\hline
\texttt{PrixParLiaison}              & Flottant  &
  Prix payé. Une valeur strictement négative trahit une annulation
  remboursée. \\
\hline
\texttt{DelaiAnticipation}           & Entier    &
  Nombre de jours entre la réservation et le départ. \\
\hline
\texttt{TypeParcoursId}              & Entier    & Type de parcours
  (national, transversal, etc.). \\
\hline
\texttt{ClassificationId}            & Entier    & Catégorie
  commerciale de la réservation. \\
\hline
\texttt{ClassePhysiqueId}            & Entier    & Classe physique
  (1ère, 2ème). \\
\hline
\texttt{NiveauPrixId}                & Entier    & Palier tarifaire. \\
\hline
\texttt{TrainAutocarId}              & Entier    & Mode de transport. \\
\hline
\texttt{CarteClientId}               & Entier    & Type de carte
  fidélité éventuelle. \\
\hline
\end{tabularx}
\caption{Principales colonnes du dataset brut et leur sémantique.}
\end{table}

\subsection{Phase d'Exploration et Identification des Problèmes}

L'exploration a été menée à l'aide de \texttt{pandas} et de scripts
ad~hoc, en parcourant systématiquement~: distribution univariée des
colonnes, valeurs manquantes, cardinalités, outliers temporels et
financiers. Les problèmes suivants ont été détectés~:

\textbf{(1) Annulations masquées.} Le dataset ne comporte aucun
\textit{flag} explicite « \textit{réservation annulée} ». Les annulations
se manifestent par~:
\begin{itemize}
\item une ligne avec \texttt{NbrVoySegment} $\leq 0$ (le segment ne
    transporte plus personne)~; ou
\item une ligne avec \texttt{PrixParLiaison} $< 0$ (un montant
    remboursé apparaît en débit).
\end{itemize}
\textbf{Implication critique~:} si l'on supprime uniquement la ligne
d'annulation, on conserve la ligne de réservation initiale, ce qui
biaise les métriques d'activité du client. Il faut supprimer
\textit{les deux} lignes corrélées (la réservation et son annulation).
Cette règle de \textbf{propagation des annulations} est appliquée dans
le module de nettoyage.

\textbf{(2) Colonnes constantes.} Plusieurs colonnes du brut ont une
cardinalité de 1 (toutes les lignes portent la même valeur). Elles
n'apportent aucune information discriminante et sont retirées
automatiquement, après détection.

\textbf{(3) Champs critiques manquants.} Certaines lignes
n'ont pas de \texttt{CodeClient}, ou pas de \texttt{LiaisonId}, ou pas
de date de départ parsable. Ces lignes sont inutilisables et rejetées.

\textbf{(4) \texttt{CodeClient} non numérique.} Quelques lignes
contiennent des valeurs parasites (chaînes non numériques) dans la
colonne \texttt{CodeClient}. Un cast en \texttt{int64} via
\texttt{pd.to\_numeric(errors="coerce")} les détecte et les exclut.

\textbf{(5) Clients \textit{cold start}.} Sur 374\,044 clients
distincts dans le dataset brut, \textbf{304\,595 (soit 81\,\%)} ont
strictement moins de 3 réservations validées après les étapes
précédentes de nettoyage. Ces clients ne peuvent pas être servis par
le modèle (par décision métier), et leurs lignes sont retirées de
l'ensemble d'entraînement. Ils représentent toutefois 403\,615 lignes
sur 491\,680 totales avant cette dernière étape, soit une part
significative du volume initial.

\textbf{(6) Couverture de la table \texttt{Liaison.csv}.} La jointure
entre les réservations et la table de référence présente un taux de
recouvrement de \textbf{99{,}92\,\%}. Les 0{,}08\,\% restants
correspondent à des identifiants de liaison absents de la table de
référence~; ces lignes sont conservées (la liaison reste un identifiant
exploitable) mais signalées dans le rapport.

\subsection{Pipeline de Nettoyage~: Algorithme}

L'ensemble des règles métier ci-dessus est implémenté dans la fonction
\texttt{make\_clean\_dataset()} du module
\texttt{src/rec\_oncf/cleaning.py}. Le pipeline applique les étapes
suivantes, dans l'ordre exact donné~:

\begin{enumerate}
\item \textbf{Vérification des colonnes requises.} La fonction lève
    immédiatement une exception si une colonne attendue est absente du
    fichier source. Cette \textit{fail-fast} évite de continuer un
    traitement avec un schéma corrompu.

\item \textbf{Parsing des dates.} La colonne
    \texttt{DateHeureDepartVoyageSegment} est convertie en
    \texttt{datetime64[ns]} via \texttt{pd.to\_datetime}. Les valeurs
    non parsables produisent un \texttt{NaT} qui sera filtré en étape~7.

\item \textbf{Suppression des colonnes constantes.} Pour chaque colonne,
    on calcule \texttt{nunique()}~; celles qui valent~1 sont écartées et
    journalisées dans le rapport de nettoyage.

\item \textbf{Normalisation de \texttt{LiaisonId}.} Les valeurs sont
    castées successivement en \texttt{Int64} puis \texttt{str}, afin
    d'uniformiser des formats mixtes (numérique/chaîne) et d'éliminer
    les zéros de tête éventuels.

\item \textbf{Jointure avec \texttt{Liaison.csv}.} Une jointure gauche
    enrichit le dataset avec les libellés de gares O/D. Le taux de
    recouvrement est calculé et reporté.

\item \textbf{Suppression des annulations.}
    Les lignes telles que \texttt{NbrVoySegment} $\leq 0$
    \textbf{ou} \texttt{PrixParLiaison} $< 0$ sont identifiées comme
    annulations. La règle de \textit{propagation} retire à la fois la
    ligne d'annulation et la réservation initiale correspondante (même
    \texttt{CodeClient}, même \texttt{LiaisonId}, fenêtre temporelle
    proche).

\item \textbf{Suppression des champs critiques manquants.}
    Toute ligne avec \texttt{CodeClient}, \texttt{LiaisonId} ou date de
    départ manquant est définitivement écartée.

\item \textbf{Déduplication.} Suppression des doublons exacts sur
    l'ensemble des colonnes après nettoyage. Cela traite les éventuels
    artefacts d'extraction (lignes en double issues d'un \texttt{JOIN}
    SQL en amont).

\item \textbf{Encodage cyclique des temporalités.} Création de
    \texttt{depart\_hour}, \texttt{depart\_dow}, \texttt{depart\_month}
    depuis la date de départ.

\item \textbf{Filtrage cold start.} Les clients ayant strictement moins
    de \texttt{MIN\_TRIPS\_FOR\_TRAINING~=~3} réservations validées
    après toutes les étapes précédentes sont retirés.
\end{enumerate}

\subsection{Résultats Quantifiés du Nettoyage}

\begin{table}[h]
\centering
\small
\begin{tabularx}{\textwidth}{|X|r|}
\hline
\textbf{Indicateur} & \textbf{Valeur} \\
\hline
Lignes brutes en entrée                       & 946\,155 \\
\hline
Lignes propres en sortie                      & 491\,680 \\
\hline
Réduction nette                               & $-48{,}0\,\%$ \\
\hline
Lignes retirées par filtre cold start         & 403\,615 \\
\hline
Clients cold start retirés                    & 304\,595 \\
\hline
Clients actifs ($\geq 3$ réservations) retenus & 69\,449 \\
\hline
Liaisons distinctes finales                   & 1\,067 \\
\hline
Taux de recouvrement avec \texttt{Liaison.csv} & 99{,}92\,\% \\
\hline
\end{tabularx}
\caption{Statistiques détaillées du nettoyage (source~:
\texttt{reports/cleaning\_report.json}).}
\end{table}

\subsection{Artefacts Produits par l'Étape de Nettoyage}

L'exécution du script \texttt{scripts/01\_make\_dataset.py} produit
trois artefacts persistés~:
\begin{itemize}
\item \texttt{data/processed/oncf\_clean.parquet}~: dataset nettoyé,
    491\,680 lignes, format Parquet pour la rapidité de lecture.
\item \texttt{reports/cleaning\_report.json}~: synthèse JSON des
    indicateurs (volumétrie avant/après, comptages d'écarts, taux de
    recouvrement).
\item \texttt{reports/cleaning\_provenance.parquet}~: trace de
    l'origine de chaque suppression (utile pour audit).
\end{itemize}

\newpage
% ================================================================
\section{Ingénierie des Variables (Feature Engineering)}
% ================================================================

\subsection{Démarche}

L'ingénierie des variables consiste à transformer le dataset propre en
une \textbf{table de features} consommable par un algorithme
d'apprentissage. Trois objectifs guident cette étape~:
\begin{enumerate}
\item \textbf{Encoder l'historique~:} les patterns de mobilité d'un
    voyageur ne se lisent pas dans une seule réservation, mais dans la
    \textit{séquence} de ses réservations. Il faut donc dériver des
    variables qui résument cette séquence.
\item \textbf{Encoder le temps~:} la prochaine liaison dépend fortement
    du moment de l'année, du jour de la semaine et de l'heure de la
    journée. L'encodage doit refléter la \textit{cyclicité} de ces
    variables.
\item \textbf{Conserver le pivot temporel~:} la date de départ doit
    rester disponible pour le split temporel ultérieur, mais ne doit
    \textbf{pas} être utilisée comme feature directement (sinon le modèle
    apprend des artefacts du futur).
\end{enumerate}

\subsection{Variables Dérivées de l'Historique}

À partir des données triées par \texttt{(CodeClient, date)},
quatre variables clés sont calculées~:

\begin{itemize}
\item \texttt{user\_trip\_index} (\textit{int}, $\geq 0$)~:
    \textit{cumulative count} du voyage dans l'historique du client.
    Un client à son premier voyage a la valeur~0, son deuxième~1, etc.
    Sémantiquement~: \textit{ancienneté du voyageur}.
\item \texttt{prev\_liaison} (\textit{string})~: la liaison immédiatement
    précédente du même client. \texttt{NaN} pour le premier voyage.
    Capture l'\textit{inertie de trajet}~: un client qui a fait
    Casablanca--Rabat hier a une probabilité élevée de refaire ce trajet
    aujourd'hui.
\item \texttt{days\_since\_prev} (\textit{float})~: nombre de jours
    écoulés depuis la réservation précédente. \texttt{NaN} pour le
    premier voyage. Capture le \textit{rythme} de l'utilisateur~:
    quotidien, hebdomadaire, occasionnel.
\item \texttt{is\_self\_purchase} (\textit{int}, 0/1)~: indicateur
    binaire valant~1 si \texttt{AchteurId == CodeClient}. Distingue
    les voyages personnels des achats pour autrui (cas typique~: parent
    qui achète pour un enfant).
\end{itemize}

\subsection{Variables Temporelles et Encodage Cyclique}

Les variables \texttt{depart\_hour} (0--23), \texttt{depart\_dow}
(0--6, 0~=~lundi) et \texttt{depart\_month} (1--12) sont par nature
\textbf{cycliques}~: l'heure 23 est proche de l'heure 0 (passage de
minuit), pas distante de 23~unités. De même, le mois 12 (décembre) est
proche du mois 1 (janvier).

Un encodage purement entier ne peut représenter cette proximité~: il
forcerait le modèle à apprendre artificiellement des règles de
\textit{cas particulier} aux frontières du cycle. La solution standard
est l'\textbf{encodage cyclique trigonométrique}, qui projette chaque
variable cyclique sur le cercle unité via une paire
$(\sin, \cos)$~:

\[
\text{sin\_h} = \sin\!\left(\frac{2\pi \cdot h}{H}\right),\qquad
\text{cos\_h} = \cos\!\left(\frac{2\pi \cdot h}{H}\right)
\]

où $h$ est la valeur de la variable et $H$ sa période ($H = 24$ pour
l'heure, $H = 7$ pour le jour de la semaine, $H = 12$ pour le mois).

Cet encodage présente deux avantages~:
\begin{enumerate}
\item Les valeurs proches sur le cycle ont des représentations
    vectorielles proches dans $\mathbb{R}^2$.
\item Les arbres de décision (utilisés par XGBoost) exploitent
    naturellement ces deux variables conjointement, en effectuant des
    splits sur l'une ou l'autre selon le signal le plus discriminant.
\end{enumerate}

Au total, six colonnes supplémentaires sont créées~:
\texttt{depart\_hour\_sin}, \texttt{depart\_hour\_cos},
\texttt{depart\_dow\_sin}, \texttt{depart\_dow\_cos},
\texttt{depart\_month\_sin}, \texttt{depart\_month\_cos}.

\subsection{Variable de Fidélité \texttt{user\_top\_liaison\_share}}

Ajoutée en Sprint~2, cette variable mesure la \textbf{fidélité} d'un
voyageur à sa ligne dominante. Elle est définie, pour le voyage à
l'instant~$t$ d'un client, comme la fraction de ses voyages
\textit{strictement antérieurs} à $t$ effectués sur sa liaison la
plus fréquente~:

\[
\text{user\_top\_liaison\_share}_t \;=\;
   \frac{\max_{\ell \in \mathcal{L}} \big|\{\,i : i < t,\; \ell_i = \ell \,\}\big|}
        {\big|\{\,i : i < t \,\}\big|}
\]

Sémantiquement, une valeur proche de~$1$ caractérise un
\textit{navetteur} (un client dont presque tous les voyages sont sur
la même ligne, typiquement domicile/travail), tandis qu'une valeur
proche de~$0{,}3$ caractérise un voyageur diversifié.
La feature est calculée \textit{en ligne} sans aucune connaissance
du voyage en cours, ce qui élimine toute fuite temporelle. Pour le
premier voyage d'un client, la valeur est \texttt{NaN} (XGBoost gère
nativement les manquants). Cette variable porte un signal complémentaire
aux variables de fréquence/récence déjà présentes (\texttt{prev\_liaison},
\texttt{user\_trip\_index}) et aide le classifieur à pondérer entre
\textit{«~refaire la ligne dominante~»} et \textit{«~prendre une ligne
secondaire~»}.

\subsection{Schéma Complet de la Table de Features}

La table \texttt{features.parquet} produite contient
\textbf{491\,680 lignes $\times$ 26 colonnes}. Ces colonnes se
répartissent en quatre catégories.

\begin{table}[h]
\centering
\small
\begin{tabularx}{\textwidth}{|l|l|X|}
\hline
\textbf{Catégorie} & \textbf{Nombre} & \textbf{Colonnes} \\
\hline
Pilotage (non features)   & 2  &
  \texttt{DateHeureDepartVoyageSegment} (split),
  \texttt{LiaisonId} (cible) \\
\hline
Lookup (jamais feature)   & 1  &
  \texttt{CodeClient} \\
\hline
Catégorielles (encodage ordinal) & 7  &
  \texttt{TypeParcoursId}, \texttt{ClassificationId},
  \texttt{ClassePhysiqueId}, \texttt{NiveauPrixId},
  \texttt{TrainAutocarId}, \texttt{CarteClientId},
  \texttt{prev\_liaison} \\
\hline
Numériques (passthrough)  & 16 &
  \texttt{PrixParLiaison}, \texttt{NbrVoySegment},
  \texttt{DelaiAnticipation}, \texttt{user\_trip\_index},
  \texttt{days\_since\_prev}, \texttt{user\_top\_liaison\_share},
  \texttt{depart\_hour}, \texttt{depart\_dow}, \texttt{depart\_month},
  \texttt{depart\_hour\_sin/cos}, \texttt{depart\_dow\_sin/cos},
  \texttt{depart\_month\_sin/cos}, \texttt{is\_self\_purchase} \\
\hline
\end{tabularx}
\caption{Répartition des 26 colonnes de la table de features.}
\end{table}

\subsection{Considérations de Robustesse Pandas~3}

Le projet est exécuté sur \texttt{pandas~3.x}, qui introduit le type
\texttt{StringDtype} comme représentation par défaut des chaînes.
Conséquence importante~: la comparaison \texttt{dtype == object} retourne
désormais \texttt{False} pour les colonnes de chaînes, contrairement à
\texttt{pandas~2}. La détection des colonnes catégorielles utilise donc
explicitement \texttt{pd.api.types.is\_numeric\_dtype()} et
\texttt{pd.api.types.is\_string\_dtype()}, garantissant la portabilité
entre les versions.

\newpage
% ================================================================
\section{Conception du Système}
% ================================================================

\subsection{Vue d'Ensemble Architecturale}

Le système se décompose en \textbf{trois couches} concentriques,
inspirées des bonnes pratiques de séparation des préoccupations~:

\begin{enumerate}
\item \textbf{Couche données et modèle} (offline)~: scripts numérotés
    \texttt{01}, \texttt{02}, \texttt{03} qui produisent
    successivement le dataset propre, la table de features et le modèle
    sérialisé.
\item \textbf{Couche métier} (\texttt{src/rec\_oncf/})~: bibliothèque
    Python importable, contenant la logique de recommandation
    (\texttt{Recommender}, \texttt{generate\_candidates}, etc.) sans
    aucune dépendance HTTP.
\item \textbf{Couche transport} (\texttt{apps/api/})~: application
    FastAPI, fine couche HTTP qui valide les requêtes et délègue à la
    couche métier.
\end{enumerate}

Cette séparation garantit que la logique métier est testable
indépendamment du transport~: les 26 tests unitaires manipulent
directement le \texttt{Recommender}, sans démarrer de serveur HTTP.

\subsection{Architecture Deux Étapes en Détail}

L'architecture deux étapes est le cœur du système.

\textbf{Étape~1~: Candidate Generation (heuristique).} Pour un
\texttt{code\_client} donné~:
\begin{enumerate}
\item Récupération de l'historique complet du client depuis le lookup
    en mémoire (clé~: \texttt{CodeClient}).
\item Conservation des 50 dernières réservations
    (\texttt{lookback=50}).
\item Comptage de la fréquence d'apparition de chaque
    \texttt{LiaisonId} dans cette fenêtre.
\item Calcul de la récence (date la plus récente d'apparition) pour
    chaque liaison.
\item Tri \textbf{décroissant par fréquence}, puis \textbf{décroissant
    par récence} (tie-break en cas d'égalité de fréquence).
\item Sélection des \textbf{10 premières liaisons}.
\end{enumerate}

Le résultat est une liste ordonnée d'au plus 10 candidats, à laquelle
le modèle de ranking sera ensuite appliqué.

\textbf{Étape~2~: Ranking (modèle XGBoost) avec filtrage par candidats.}
\begin{enumerate}
\item Récupération de la \textit{dernière ligne de features} du client
    (clé~: \texttt{CodeClient}, source~:
    \texttt{features.parquet} en mémoire).
\item Application du \texttt{ColumnTransformer} entraîné~:
    \texttt{OrdinalEncoder} sur les 7 colonnes catégorielles, passthrough
    sur les 16 numériques.
\item Inférence du modèle~: \texttt{predict\_proba(X)} retourne un
    vecteur de probabilités sur les 1\,011 classes.
\item \textbf{Filtrage par candidats~:} les probabilités sont
    \textit{restreintes aux indices correspondant aux candidats produits
    par l'étape~1}, via \texttt{label\_encoder.transform(candidates)}.
    Les candidats inconnus du modèle (cas marginal) sont écartés ; si
    aucun candidat n'est connu, la sortie heuristique brute de
    l'étape~1 est utilisée comme \textit{fallback}.
\item Tri décroissant \textit{parmi les candidats restreints} par leur
    score, sélection des $k$ premiers ($k \in \{1, 2, 3\}$).
\item Renvoi des \texttt{LiaisonId} correspondants.
\end{enumerate}

C'est cette étape de filtrage explicite qui réalise effectivement
l'architecture deux étapes décrite plus haut~: sans elle, le top-$k$
serait pris parmi toutes les classes du modèle et l'étape~1 ne
servirait à rien.

\subsection{Règle de Cold Start}

Le système retourne explicitement un mode \texttt{cold\_start} dans
trois situations~:
\begin{enumerate}
\item Le \texttt{code\_client} fourni n'existe pas dans le lookup
    historique.
\item L'historique du client contient \textit{strictement moins de}
    3~réservations.
\item La \texttt{generate\_candidates()} retourne une liste vide
    (cas pathologique d'historique non vide mais sans liaison
    exploitable).
\end{enumerate}

Dans ces trois cas, la réponse JSON est~:
\texttt{\{"mode": "cold\_start", "recommendations": []\}}.

Le choix de retourner explicitement un mode --- plutôt que de tenter une
recommandation par défaut --- a été motivé par le souci de transparence
métier~: l'application mobile peut alors décider d'afficher un message
contextuel à l'utilisateur (\textit{« nous apprenons encore vos
habitudes »}) plutôt qu'une suggestion potentiellement non pertinente.

\subsection{Lookups en Mémoire au Démarrage}

À l'initialisation du \texttt{Recommender} (méthode
\texttt{from\_paths}), deux structures de données sont construites en
mémoire~:
\begin{itemize}
\item \texttt{history\_lookup~: dict[str, DataFrame]}~: pour chaque
    \texttt{CodeClient}, l'intégralité de son historique trié par date.
    Source~: \texttt{oncf\_clean.parquet}.
\item \texttt{feature\_lookup~: dict[str, DataFrame]}~: pour chaque
    \texttt{CodeClient}, sa dernière ligne de features (la plus récente).
    Source~: \texttt{features.parquet}.
\end{itemize}

La construction est effectuée une fois pour toutes au démarrage du
processus~: les requêtes ultérieures sont des accès $O(1)$ aux
dictionnaires en mémoire. Cette stratégie sacrifie une portion de RAM
(estimée à $\approx 350$~Mo, en deçà du budget SNF-03) pour garantir
une latence négligeable.

\subsection{Diagramme de Séquence d'une Requête}

Le flux d'une requête \texttt{POST /recommend} est le suivant~:

\begin{enumerate}
\item Le client mobile envoie une requête HTTP avec le corps JSON
    \texttt{\{"code\_client": "12345", "k": 3\}}.
\item FastAPI parse et valide le corps via Pydantic ($k$ contraint à
    $[1, 3]$).
\item Le \texttt{Recommender.recommend()} est appelé avec les
    paramètres validés.
\item Lookup historique~: si absent ou $< 3$ lignes
    $\rightarrow$ retour \texttt{cold\_start}.
\item Sinon~: \texttt{generate\_candidates(history)}.
\item Lookup features~: récupération de la dernière ligne pour
    construire la matrice $X$ d'inférence.
\item Inférence XGBoost~: \texttt{predict\_proba(X)} sur les 1\,011
    classes.
\item Sélection top-$k$~: tri décroissant et sélection.
\item Décodage des indices vers \texttt{LiaisonId}.
\item Construction et renvoi de la réponse JSON
    \texttt{\{"mode": "model", "recommendations": [...]\}}.
\end{enumerate}

\newpage
% ================================================================
\section{Modélisation~: XGBoost Multiclasse}
% ================================================================

\subsection{Formulation du Problème}

Le problème de prédiction est formulé comme une \textbf{classification
multiclasse} sur 1\,011 classes (le nombre de \texttt{LiaisonId}
distincts dans le jeu d'entraînement après split temporel).

Pour une ligne de features $\mathbf{x} \in \mathbb{R}^p$ (avec $p =
22$ après préprocessing), le modèle prédit un vecteur de probabilités
$\hat{\mathbf{p}} \in [0,1]^{1011}$ tel que $\sum_{k=0}^{1010}
\hat{p}_k = 1$. La classe prédite est~:

\[
\hat{c}(\mathbf{x}) = \arg\max_{k \in \{0, \ldots, 1010\}} \hat{p}_k(\mathbf{x})
\]

\subsection{Pourquoi XGBoost~?}

Plusieurs algorithmes ont été considérés en début de modélisation~:
\begin{itemize}
\item \textbf{Régression logistique multinomiale}~: trop simple, ne
    capture pas les interactions entre features.
\item \textbf{Random Forest}~: solide, mais moins performant que
    XGBoost sur les problèmes tabulaires de cette taille.
\item \textbf{Réseaux de neurones tabulaires (TabNet, FT-Transformer)}~:
    intéressants, mais sur-dimensionnés pour 22 features et
    491\,680 lignes~; la complexité d'entraînement (tuning, GPU,
    saveur \textit{deep learning}) n'est pas justifiée.
\item \textbf{XGBoost}~: standard de l'industrie sur les tâches
    tabulaires, gère nativement le multiclasse, propose une
    régularisation efficace, est rapide à entraîner et produit des
    feature importances exploitables.
\end{itemize}

XGBoost (\textit{eXtreme Gradient Boosting}) est une implémentation
optimisée de \textit{Gradient Boosted Decision Trees}, introduite par
Chen et Guestrin (2016). Elle combine~: arbres CART de régression
empilés, descente de gradient sur la fonction objectif, régularisation
L1~/~L2, échantillonnage des lignes et des colonnes,
parallélisation~CPU/GPU.

\subsection{Prétraitement~: \texttt{ColumnTransformer}}

Un \texttt{ColumnTransformer} sklearn est placé en amont du
classifieur. Il applique~:
\begin{itemize}
\item Sur les \textbf{7 colonnes catégorielles} (\texttt{TypeParcoursId},
    \texttt{ClassificationId}, \texttt{ClassePhysiqueId},
    \texttt{NiveauPrixId}, \texttt{TrainAutocarId}, \texttt{CarteClientId},
    \texttt{prev\_liaison})~: un \texttt{OrdinalEncoder} avec
    \texttt{handle\_unknown="use\_encoded\_value"} et
    \texttt{unknown\_value=-1}. Chaque modalité reçoit un entier
    ordinal~; les modalités inconnues à l'inférence sont mappées vers
    $-1$ (et non rejetées, ce qui briserait la robustesse en
    production).
\item Sur les \textbf{16 colonnes numériques}~: \texttt{passthrough}
    (aucune transformation, les arbres sont insensibles à l'échelle).
\end{itemize}

\subsubsection*{Pourquoi OrdinalEncoder et non OneHotEncoder~?}

La colonne \texttt{prev\_liaison} possède \textbf{1\,011 valeurs
distinctes}. Un \texttt{OneHotEncoder} générerait une matrice de
features de plus de \textbf{5\,000 colonnes} (1\,011 pour
\texttt{prev\_liaison} + quelques dizaines pour les autres
catégorielles). Cela aurait trois conséquences délétères~:
\begin{enumerate}
\item Empreinte mémoire multipliée par $\approx 200$.
\item Temps d'entraînement et d'inférence considérablement allongés.
\item Sparsité élevée, peu favorable aux splits des arbres.
\end{enumerate}

L'\texttt{OrdinalEncoder} maintient la matrice à $\approx 23$ colonnes
denses, exploitables directement par XGBoost. Les arbres de décision
gèrent nativement les variables ordinales~: ils n'imposent pas d'ordre
naturel mais apprennent des partitions arbitraires de l'espace de
modalités.

\subsection{Encodage du Label}

La variable cible \texttt{LiaisonId} est de type chaîne. XGBoost attend
des labels entiers $\{0, \ldots, K-1\}$. Un
\texttt{LabelEncoder} sklearn opère la conversion~:
\[
\texttt{"4512"} \xrightarrow{\text{LabelEncoder}} 42
\]

L'encodeur est sérialisé avec le modèle (artefact
\texttt{label\_encoder.joblib}) afin de permettre le décodage
\texttt{inverse\_transform()} à l'inférence.

\subsection{Hyperparamètres Retenus}

Les hyperparamètres finaux ont été ajustés après plusieurs itérations,
en arbitrant entre qualité de prédiction et faisabilité (saturation
mémoire CPU).

\begin{table}[h]
\centering
\small
\begin{tabularx}{\textwidth}{|l|c|X|}
\hline
\textbf{Paramètre}      & \textbf{Valeur}        & \textbf{Justification} \\
\hline
\texttt{objective}      & \texttt{multi:softprob} & Sortie probabiliste sur 1\,011 classes \\
\hline
\texttt{eval\_metric}   & \texttt{mlogloss}       & Log-vraisemblance multiclasse, cohérente avec l'objectif \\
\hline
\texttt{tree\_method}   & \texttt{hist}           & Algorithme histogram-based, rapide et économe en mémoire \\
\hline
\texttt{device}         & \texttt{cpu}            & Réduit à CPU après échec OOM sur GPU (cf. section~9.1) \\
\hline
\texttt{n\_estimators}  & 200                     & Réduit de 300 pour stabilité CPU \\
\hline
\texttt{learning\_rate} & 0.08                    & Compromis convergence/généralisation \\
\hline
\texttt{max\_depth}     & 6                       & Réduit de 8 pour stabilité CPU et lutte contre l'overfitting \\
\hline
\texttt{subsample}      & 0.9                     & Échantillonnage 90\,\% des lignes par arbre, régularisation \\
\hline
\texttt{colsample\_bytree} & 0.8                  & Échantillonnage 80\,\% des colonnes par arbre, régularisation \\
\hline
\texttt{reg\_lambda}    & 1.0                     & Régularisation L2 sur les poids des feuilles \\
\hline
\texttt{n\_jobs}        & $-1$                    & Utilisation de tous les cœurs CPU disponibles \\
\hline
\end{tabularx}
\caption{Hyperparamètres retenus pour le modèle XGBoost.}
\end{table}

\subsection{Fonction Objectif et Régularisation}

L'objectif \texttt{multi:softprob} optimise la
\textbf{log-vraisemblance multiclasse} (mlogloss)~:

\[
\mathcal{L}_{\text{data}} = -\sum_{i=1}^{N} \log P(y_i \mid \mathbf{x}_i)
\]

où $P(y_i \mid \mathbf{x}_i)$ est la probabilité prédite par la
softmax~:

\[
P(y = k \mid \mathbf{x}) =
  \frac{\exp(s_k(\mathbf{x}))}
       {\displaystyle\sum_{j=0}^{1010} \exp(s_j(\mathbf{x}))}
\]

et $s_k(\mathbf{x})$ est le score brut (\textit{logit}) de la classe $k$,
calculé comme la somme des contributions des arbres successifs.

À cette perte de données s'ajoute un terme de régularisation
$\Omega(f)$ pour chaque arbre $f$~:
\[
\Omega(f) = \gamma T + \tfrac{1}{2}\lambda \|w\|^2
\]
où $T$ est le nombre de feuilles de l'arbre, $w$ le vecteur des poids
des feuilles, $\gamma$ pénalise la complexité structurelle (nombre de
feuilles) et $\lambda$ pénalise l'amplitude des poids.
La fonction objectif globale optimisée est~:

\[
\mathcal{O} = \mathcal{L}_{\text{data}} + \sum_{f \in F} \Omega(f)
\]

XGBoost minimise $\mathcal{O}$ par approximation de second ordre de
$\mathcal{L}_{\text{data}}$ et construction gloutonne d'arbres
maximisant le \textit{gain} à chaque split.

\subsection{Méthodologie de Validation~: Split Temporel}

Le choix du protocole de validation est crucial. Un \textbf{split
aléatoire} (\textit{train\_test\_split} classique) permettrait au
modèle d'« apprendre » des voyages futurs pour prédire des voyages
passés~: il y aurait une fuite d'information temporelle qui
surestimerait artificiellement les métriques.

Le \textbf{split temporel} simule au contraire les conditions réelles
de production~:
\begin{enumerate}
\item Le dataset est trié par \texttt{DateHeureDepartVoyageSegment}
    croissant.
\item Les \textbf{80\,\%} premiers voyages constituent le jeu
    d'entraînement (393\,344 lignes).
\item Les \textbf{20\,\%} restants constituent le jeu de test
    (98\,261 lignes après filtrage des 75 labels absents du train).
\end{enumerate}

Cette discipline garantit que les métriques rapportées reflètent
fidèlement la performance attendue en production.

\newpage
% ================================================================
\section{Difficultés Rencontrées et Résolutions}
% ================================================================

Quatre difficultés significatives ont jalonné le développement du
modèle. Leur description ici remplit deux rôles~: (i)~documenter pour
la reproductibilité, et (ii)~expliciter le raisonnement technique qui a
abouti aux choix actuels.

\subsection{Difficulté~1~: Saturation VRAM GPU et OOM Silencieux}

\textbf{Symptôme observé.} Les premières tentatives d'entraînement
utilisaient \texttt{device="cuda"} pour exploiter le GPU NVIDIA RTX~3050
disponible sur la machine. À environ 40\,\% de l'avancée
(arbres~80--120 sur 300), le processus Python disparaissait
\textit{silencieusement}~: aucune trace de pile, aucun message
d'erreur, simplement le retour brutal au prompt.

\textbf{Diagnostic.} L'analyse a montré que le RTX~3050 ne dispose que
de \textbf{4~Go de VRAM}. Or, XGBoost en mode CUDA réserve une
matrice de gradients de taille proportionnelle à
$N \times K$ où $N = 393\,344$ et $K = 1\,011$. Couplé aux
structures internes des arbres (\texttt{max\_depth = 8}) et au nombre
d'estimateurs (300), le total dépassait la VRAM disponible. Le \textit{OOM
Killer} du système d'exploitation (Windows en l'occurrence) terminait le
processus sans propager d'exception au niveau Python.

\textbf{Résolution.} Trois leviers ont été combinés~:
\begin{enumerate}
\item Bascule vers \texttt{device="cpu"} (le plus important).
\item Réduction de \texttt{n\_estimators} de 300 à \textbf{200}.
\item Réduction de \texttt{max\_depth} de 8 à \textbf{6}.
\end{enumerate}

Le temps d'entraînement passe à $\approx 43$~minutes sur CPU. Les
métriques obtenues restent largement au-dessus des seuils, ce qui
montre que le couple \texttt{n\_estimators=200, max\_depth=6} est un
\textit{sweet spot} pour ce dataset.

\subsection{Difficulté~2~: Fuite du \texttt{CodeClient} comme Feature}

\textbf{Symptôme observé.} Une revue de code a révélé que, dans une
version intermédiaire de \texttt{train\_xgb\_multiclass}, la matrice
$X$ utilisée pour l'entraînement n'excluait pas explicitement la
colonne \texttt{CodeClient}. Le modèle l'utilisait donc comme variable
prédictive.

\textbf{Pourquoi c'est grave.} D'une part, c'est une violation directe
des engagements de conformité Loi~09-08~: l'identifiant client devient
implicitement une variable du modèle, ce qui rend ce dernier
\textit{auditable} comme traitant l'identifiant comme une donnée
exploitée. D'autre part, l'identifiant numérique est arbitraire~:
le modèle peut sur-apprendre des associations spurieuses
(client~12345 $\rightarrow$ liaison~42) qui ne se généralisent pas, et
qui dégradent la qualité réelle au lieu de l'améliorer.

\textbf{Résolution.} Suppression explicite et défensive dans
\texttt{train\_xgb\_multiclass} et \texttt{predict\_proba}~:

\begin{lstlisting}[style=py]
_drop = [c for c in [label_col, time_col, "CodeClient"]
         if c in df_train.columns]
X = df_train.drop(columns=_drop)
\end{lstlisting}

Cette construction défensive (filtrage avec \texttt{if c in
df\_train.columns}) tolère l'absence éventuelle d'une colonne sans
lever d'exception. Elle est doublée d'un test unitaire dédié dans
\texttt{tests/test\_training.py}.

\subsection{Difficulté~3~: Labels Inconnus à l'Évaluation}

\textbf{Symptôme observé.} Lors de la première exécution du
script~03, une exception \texttt{ValueError} était levée par
\texttt{label\_encoder.transform()} appliqué au jeu de test.

\textbf{Diagnostic.} Sur les 1\,067 \texttt{LiaisonId} distincts
présents dans le dataset propre, seules \textbf{1\,011} apparaissent
dans la portion d'entraînement (les 80\,\% temporellement
antérieurs). Les \textbf{56 restantes} apparaissent uniquement dans la
portion de test --- ce sont des liaisons \textit{nouvelles} qui n'ont
émergé que dans la période la plus récente. Comme le
\texttt{LabelEncoder} est entraîné sur le seul jeu d'entraînement, il
ignore ces classes nouvelles. La transformation crash en conséquence.

Une analyse fine montre que cela représente
\textbf{75 lignes du jeu de test} (sur 98\,336 initialement).

\textbf{Résolution.} Filtrage du jeu de test avant évaluation, pour
ne conserver que les lignes dont la classe est connue de l'encodeur~:

\begin{lstlisting}[style=py]
known_classes = set(artifacts.label_encoder.classes_)
test_df = test_df[
    test_df[label_col].astype(str).isin(known_classes)
]
\end{lstlisting}

Le rapport de métriques journalise explicitement le nombre de lignes
retirées sous le champ \texttt{rows\_test\_unseen\_dropped} (75), pour
auditabilité.

\subsection{Difficulté~4~: Incohérence d'Encodeur entre Tests et Production}

\textbf{Symptôme observé.} Les tests unitaires de
\texttt{tests/test\_recommender.py} échouaient en signalant une
inadéquation de dimensions entre la matrice de features attendue par
le modèle et celle produite par le \texttt{ColumnTransformer}.

\textbf{Diagnostic.} Les fixtures de test utilisaient initialement un
\texttt{OneHotEncoder} pour simplifier l'écriture du test. Or, la
production utilise un \texttt{OrdinalEncoder}~; la dimension de sortie
du \texttt{ColumnTransformer} diffère drastiquement (sparse $\sim 5000$
colonnes pour OHE vs dense $\sim 23$ pour OE). Le modèle entraîné en
production attend la dimension dense.

\textbf{Résolution.} Remplacement du \texttt{OneHotEncoder} par
\texttt{OrdinalEncoder} dans toutes les fixtures de test, alignant
ainsi la pile de test avec la pile de production. Cela évite des
défaillances en intégration.

\textbf{Leçon générale.} Les fixtures de test doivent répliquer
fidèlement le pipeline de production, même au prix d'une légère
complexité supplémentaire. Une simplification illusoire en test
masque la véritable surface de bug.

\newpage
% ================================================================
\section{Méthodologie d'Évaluation et Résultats}
% ================================================================

\subsection{Choix des Métriques}

Les systèmes de recommandation se mesurent par des métriques
spécifiques qui reflètent l'expérience utilisateur réelle~: il ne
suffit pas que le modèle « ait raison souvent », il faut que la bonne
réponse apparaisse \textbf{en haut de la liste} présentée à
l'utilisateur. Deux familles de métriques sont utilisées dans ce
projet.

\subsubsection*{Hit Rate at $k$ (HR@$k$)}

Le \textbf{Hit Rate at $k$} mesure la proportion de cas où la vraie
prochaine liaison figure dans les $k$ premières recommandations
retournées~:
\[
\text{HR@}k = \frac{1}{N} \sum_{i=1}^{N}
  \mathbf{1}\!\left[\,\text{rang}(y_i) \leq k\,\right]
\]
où $\text{rang}(y_i)$ est le rang de la vraie classe $y_i$ dans la
liste de prédictions ordonnées par probabilité décroissante, et
$\mathbf{1}[\cdot]$ est la fonction indicatrice.

Deux variantes sont rapportées~:
\begin{itemize}
\item \textbf{HR@1}~: la première recommandation est-elle correcte~?
    C'est la métrique la plus stricte~; elle reflète l'expérience
    utilisateur la plus exigeante (\textit{« le système devine du premier
    coup »}).
\item \textbf{HR@3}~: l'une des trois premières recommandations
    est-elle correcte~? C'est la métrique pertinente pour l'écran
    d'accueil mobile, qui peut afficher trois suggestions.
\end{itemize}

\subsubsection*{Mean Reciprocal Rank at $k$ (MRR@$k$)}

Le \textbf{Mean Reciprocal Rank at $k$} pondère le succès par le rang
auquel la bonne réponse apparaît~:
\[
\text{MRR@}k = \frac{1}{N} \sum_{i=1}^{N}
  \frac{1}{\text{rang}(y_i)}
  \cdot \mathbf{1}\!\left[\,\text{rang}(y_i) \leq k\,\right]
\]

Concrètement~: un hit en position~1 vaut 1{,}0~; en position~2,
0{,}50~; en position~3, $\approx 0{,}33$. Cette métrique pénalise les
recommandations correctes mais mal classées et complète utilement le
HR.

\subsection{Implémentation des Métriques}

L'implémentation, contenue dans \texttt{src/rec\_oncf/metrics.py},
exploite le fait que le tableau \texttt{proba} de sortie de
\texttt{predict\_proba} est de forme $(N, K)$~:

\begin{lstlisting}[style=py, caption={Implémentation des métriques}]
import numpy as np

def hit_rate_at_k(y_true, proba, *, k):
    topk = np.argsort(-proba, axis=1)[:, :k]
    hits = (topk == y_true[:, None]).any(axis=1)
    return float(hits.mean())

def mrr_at_k(y_true, proba, *, k):
    topk = np.argsort(-proba, axis=1)[:, :k]
    mrr = 0.0
    for rank in range(k):
        hit = topk[:, rank] == y_true
        mrr += hit.sum() / (rank + 1)
    return float(mrr / len(y_true))
\end{lstlisting}

L'opération \texttt{np.argsort(-proba, axis=1)} retourne les indices
des classes triés par probabilité \textit{décroissante} (le signe
négatif inverse l'ordre par défaut croissant). Le slicing
\texttt{[:, :k]} ne conserve que les $k$ premiers rangs.

\subsection{Résultats Quantitatifs}

\begin{table}[h]
\centering
\small
\begin{tabularx}{\textwidth}{|X|c|c|c|c|}
\hline
\textbf{Métrique} & \textbf{Sprint~1} & \textbf{Sprint~2} & \textbf{Seuil cible} & \textbf{Validé} \\
\hline
Hit Rate@1 & 0{,}7395 & \textbf{0{,}7628} & $> 0{,}30$ & Oui (gap +154\,\%) \\
\hline
Hit Rate@3 & 0{,}8877 & \textbf{0{,}9055} & $> 0{,}50$ & Oui (gap +81\,\%) \\
\hline
MRR@3      & 0{,}8064 & \textbf{0{,}8277} & $> 0{,}35$ & Oui (gap +136\,\%) \\
\hline
\multicolumn{5}{|l|}{\small Lignes entraînement~: 393\,344 ; lignes test~: 98\,261 ; classes~: 1\,011} \\
\hline
\multicolumn{5}{|l|}{\small Lignes test retirées (labels inconnus)~: 75 ; durée d'entraînement~: $\approx 43$ minutes (CPU) ; \texttt{random\_state}~: 42} \\
\hline
\end{tabularx}
\caption{Métriques offline mesurées sur le split temporel,
avant et après ajout de la feature \texttt{user\_top\_liaison\_share}
en Sprint~2 (\texttt{reports/offline\_metrics.json}).}
\end{table}

\subsection{Métriques par Segment d'Historique Utilisateur (Sprint~2)}

L'ajout en Sprint~2 d'une métrique \textit{segmentée par profondeur
d'historique utilisateur} permet de vérifier que les performances ne
sont pas un artefact de la portion la plus facile du jeu de test. Les
segments sont définis par le nombre de réservations qu'un utilisateur
possède dans le \textit{jeu d'entraînement}.

\begin{table}[h]
\centering
\small
\begin{tabularx}{\textwidth}{|X|c|c|c|c|}
\hline
\textbf{Segment} & \textbf{$n$ test} & \textbf{HR@1} & \textbf{HR@3} & \textbf{MRR@3} \\
\hline
0--2 voyages dans train (très peu / nouveaux) & 44\,397 & 0{,}7393 & 0{,}8930 & 0{,}8091 \\
\hline
3--5 voyages (nouveaux actifs)               & 16\,780 & 0{,}7517 & 0{,}9038 & 0{,}8210 \\
\hline
6--20 voyages (réguliers)                    & 23\,737 & 0{,}7786 & 0{,}9159 & 0{,}8411 \\
\hline
21+ voyages (fidèles, navetteurs)            & 13\,347 & \textbf{0{,}8268} & \textbf{0{,}9311} & \textbf{0{,}8741} \\
\hline
\end{tabularx}
\caption{Métriques par profondeur d'historique utilisateur. La
performance s'améliore monotoniquement avec le volume d'historique~--
résultat conforme aux attentes théoriques.}
\end{table}

L'observation cruciale est la \textbf{progression monotone} des
métriques avec la profondeur d'historique~: HR@1 passe de 73{,}9\,\%
pour les utilisateurs avec moins de 3 voyages dans le train à
\textbf{82{,}7\,\%} pour les voyageurs fidèles ($\geq 21$ voyages). Cette
progression valide deux hypothèses~:

\begin{enumerate}
\item Le modèle apprend effectivement quelque chose de
    \textit{spécifique à l'utilisateur} (sinon les segments donneraient
    des performances similaires).
\item La règle \textit{cold start} ($< 3$ réservations) reste
    pertinente~: même si la performance sur les segments faibles est
    déjà élevée, les voyageurs très récents bénéficieraient d'une
    expérience utilisateur encore plus fiable avec davantage de
    données.
\end{enumerate}

\subsection{Interprétation Qualitative}

\textbf{Lecture du HR@1.} Sur 98\,261 lignes de test, le modèle prédit
\textit{exactement} la prochaine liaison comme première recommandation
dans 76{,}3\,\% des cas. Pour mettre cette valeur en perspective~:
le tirage aléatoire uniforme sur 1\,011 classes donnerait une HR@1 de
$1/1011 \approx 0{,}1\,\%$. Notre modèle est donc \textbf{770~fois plus
précis} qu'un tirage aléatoire.

\textbf{Lecture du HR@3.} En présentant les trois meilleures
suggestions sur l'écran d'accueil mobile, l'application couvrirait
\textit{nominalement} 90{,}6\,\% des prochains voyages réels~: dans plus
de neuf cas sur dix, l'utilisateur n'a qu'à toucher l'une des trois
cartes affichées plutôt que de saisir manuellement.

\textbf{Lecture du MRR@3.} La valeur de 0{,}8277 implique que, en
moyenne, la vraie liaison est classée à une position $\approx 1{,}21$
parmi les recommandations. Cela signifie qu'au-delà de la simple
présence dans le top-3, la bonne réponse est \textbf{très souvent} en
première position.

\subsection{Comparaison avec des Baselines de Référence}

Pour contextualiser les performances de XGBoost, trois baselines
classiques ont été calculées sur le \textbf{même split temporel}
(script \texttt{scripts/04\_baselines.py}, sortie
\texttt{reports/baseline\_metrics.json})~:

\begin{itemize}
\item \texttt{global\_top}~: prédire les 3~liaisons les plus populaires
    globalement (popularité absolue, indépendante de l'utilisateur).
\item \texttt{prev\_liaison}~: prédire la dernière liaison vue de
    l'utilisateur (inertie pure).
\item \texttt{most\_frequent}~: prédire les 3~liaisons les plus
    fréquentes dans l'historique de l'utilisateur (avec tie-break par
    récence, padded par popularité globale pour atteindre $k=3$).
\end{itemize}

\begin{table}[h]
\centering
\small
\begin{tabularx}{\textwidth}{|X|c|c|c|}
\hline
\textbf{Modèle}                       & \textbf{HR@1} & \textbf{HR@3} & \textbf{MRR@3} \\
\hline
\texttt{global\_top} (popularité)     & 0{,}0399      & 0{,}1125      & 0{,}0707 \\
\hline
\texttt{prev\_liaison} (dernier vu)   & 0{,}2620      & 0{,}3204      & 0{,}2881 \\
\hline
\texttt{most\_frequent} (fréquence utilisateur) & 0{,}2751 & 0{,}5128 & 0{,}3865 \\
\hline
\textbf{XGBoost multiclasse (Sprint~2)} & \textbf{0{,}7628} & \textbf{0{,}9055} & \textbf{0{,}8277} \\
\hline
\end{tabularx}
\caption{Comparaison de XGBoost avec trois baselines de référence
sur le même split temporel.}
\end{table}

\textbf{Lecture~:} XGBoost fait \textbf{2{,}77 fois mieux} que la meilleure
baseline (\texttt{most\_frequent}) en HR@1, et \textbf{1{,}77 fois mieux}
en HR@3. Le gain n'est donc pas trivial~: si l'on s'arrêtait à
«~prédire le trajet le plus fréquent du client~», on resterait à
27\,\% de réussite au top-1, ce qui ne justifierait pas une
expérience proactive. Le passage à un modèle XGBoost capable
d'exploiter le contexte (jour, heure, mois, prix, classe, type de
parcours, fidélité \texttt{user\_top\_liaison\_share}) fait passer le
HR@1 à 76\,\%, ce qui change qualitativement la valeur produit.

\textbf{Note méthodologique~:} les baselines ont été conçues comme
des heuristiques honnêtes (pas de leakage temporel, jamais
entraînées sur le test) et sont \textit{déterministes}. Elles
constituent une borne basse réaliste de ce que l'on aurait obtenu
sans modèle ML.

\subsection{Discussion sur la Sur-Performance}

Les métriques obtenues dépassent largement les seuils initialement
fixés (gap de +147\,\% pour HR@1). Plusieurs facteurs expliquent cette
sur-performance~:
\begin{enumerate}
\item \textbf{Régularité des comportements ferroviaires.} Les
    voyageurs ONCF actifs sont en grande majorité des navetteurs
    domicile/travail. Leurs habitudes sont stables et donc
    \textit{intrinsèquement prédictibles}.
\item \textbf{Filtrage cold start strict.} En écartant les clients
    avec moins de 3 réservations, on retire la portion la plus
    difficile à prédire. Le modèle est entraîné sur la portion la plus
    « régulière » du dataset.
\item \textbf{Variable \texttt{prev\_liaison}.} Cette feature à elle
    seule capte une grande partie de la variance. L'inertie de trajet
    (refaire la même liaison qu'hier) est dominante.
\item \textbf{Encodage cyclique.} Permet à XGBoost d'exploiter
    finement la périodicité temporelle (week-end vs jours ouvrés,
    matin vs après-midi).
\end{enumerate}

Une vigilance s'impose néanmoins pour les versions futures~: les
métriques sont \textit{offline} sur historique. Le passage en
production peut faire émerger des écarts (\textit{data drift},
\textit{concept drift}) qui justifieront un \textit{monitoring} en
ligne et un réentraînement périodique.

\newpage
% ================================================================
\section{Composant \texttt{Recommender} et API REST}
% ================================================================

\subsection{Le Composant \texttt{Recommender} (Logique Métier)}

Toute la logique de recommandation est encapsulée dans la classe
\texttt{Recommender} du module \texttt{src/rec\_oncf/recommender.py}.
L'API REST n'est qu'une couche fine de transport HTTP qui délègue
entièrement à cette classe~; ce découplage est volontaire (cf.
SNF-06).

\textbf{Constructeurs disponibles~:}
\begin{itemize}
\item \texttt{Recommender.from\_paths(paths)}~: charge depuis disque
    le modèle (\texttt{xgb\_ranker.json}), le \texttt{LabelEncoder}
    (\texttt{label\_encoder.joblib}) et construit en mémoire le
    \texttt{history\_lookup} à partir du Parquet propre.
\item \texttt{Recommender.from\_data(arts, clean\_df)}~:
    constructeur alternatif pour les tests unitaires, qui prend
    directement en entrée des objets en mémoire (sans passer par le
    disque). Indispensable pour des tests rapides et déterministes.
\end{itemize}

\textbf{Méthode publique principale~:}
\texttt{recommend(code\_client: str, k: int = 1, asof=None) -> dict}~:
retourne un dictionnaire avec les clés \texttt{"mode"} et
\texttt{"recommendations"}. Le paramètre \texttt{asof} (par défaut~:
heure courante) permet de simuler un instant de prédiction donné
(utile pour les tests et le \textit{backtesting}).

\subsection{Algorithme de \texttt{recommend()} (architecture deux étapes)}

\begin{lstlisting}[style=py, caption={Pseudo-code de \texttt{Recommender.recommend()}}]
def recommend(self, code_client, k=1, *, asof=None):
    history = self.history_lookup.get(code_client)
    if history is None or len(history) < 3:
        return _COLD_START

    # Etape 1 : Candidate Generation (heuristique, top 10)
    candidates = generate_candidates(history)
    if not candidates:
        return _COLD_START

    # Recalcul des features ON-THE-FLY depuis l'historique frais
    feature_row = compute_inference_row(history, asof=asof)

    # Etape 2 : Ranking XGBoost, RESTREINT aux candidats de l'etape 1
    le = self.artifacts.label_encoder
    valid = [c for c in candidates if c in set(le.classes_)]
    proba = predict_proba(self.artifacts, feature_row)[0]
    cand_scores = proba[le.transform(valid)]
    order = np.argsort(-cand_scores)[:k]
    return {"mode": "model", "recommendations": [valid[i] for i in order]}
\end{lstlisting}

\subsubsection*{Évolution Sprint~2~: features recalculées à la volée}

Dans la version initiale, le \texttt{Recommender} maintenait un second
lookup \texttt{feature\_lookup} construit au démarrage à partir du
\texttt{features.parquet} \textit{généré au moment de l'entraînement}.
Cette approche présente un défaut majeur en production~: dès qu'un
client effectue de nouvelles réservations entre deux réentraînements,
les features stockées (\texttt{prev\_liaison}, \texttt{user\_trip\_index},
\texttt{days\_since\_prev}, \texttt{user\_top\_liaison\_share}) deviennent
\textit{périmées}~; le modèle prédit à partir d'une représentation
obsolète du voyageur.

La fonction \texttt{compute\_inference\_row(history, asof=...)} introduite
en Sprint~2 résout ce problème~: elle calcule la ligne de features
\textit{à la volée}, à partir de l'historique \textit{vivant} du client,
au moment précis de la requête. Les features dérivées de l'historique
reflètent ainsi l'état réel du voyageur, et les features temporelles
cycliques sont calées sur le moment courant (\texttt{asof = now()}). Les
features de contexte de voyage (prix, classe, etc.) du futur trajet ---
inconnues par construction --- sont approchées par celles de la dernière
réservation observée.

L'usage de \texttt{dict.get(...)} (et non d'accès direct
\texttt{dict[...]}) garantit que les clés inconnues retournent
\texttt{None} sans lever \texttt{KeyError}~: c'est la sémantique
attendue pour un client jamais vu.

\subsection{Application FastAPI}

L'API REST est implémentée dans \texttt{apps/api/main.py}. Le choix de
\textbf{FastAPI} est motivé par~:
\begin{itemize}
\item Validation automatique des requêtes via \texttt{Pydantic} (v2),
    avec génération d'erreurs HTTP 422 claires en cas de payload
    invalide.
\item Documentation OpenAPI / Swagger générée automatiquement
    (\texttt{/docs}).
\item Performance native asynchrone (Starlette/Uvicorn).
\item Mécanisme de \texttt{lifespan} pour exécuter du code au
    démarrage et à l'arrêt du serveur.
\end{itemize}

\textbf{Mécanisme de chargement (\texttt{lifespan}).}
Le modèle et les lookups sont chargés \textit{une seule fois} au
démarrage du serveur, puis maintenus en mémoire pour la durée de vie
du processus~:

\begin{lstlisting}[style=py, caption={Démarrage de l'application FastAPI}]
from contextlib import asynccontextmanager
from fastapi import FastAPI
from rec_oncf.recommender import Recommender
from rec_oncf.config import default_paths

recommender: Recommender | None = None

@asynccontextmanager
async def lifespan(app: FastAPI):
    global recommender
    recommender = Recommender.from_paths(default_paths())
    yield
    # nettoyage eventuel a l'arret

app = FastAPI(lifespan=lifespan)
\end{lstlisting}

\subsection{Endpoints Exposés}

\begin{table}[h]
\centering
\small
\begin{tabularx}{\textwidth}{|l|l|X|}
\hline
\textbf{Méthode} & \textbf{Route}      & \textbf{Description} \\
\hline
GET  & \texttt{/health}    &
  Probe de santé pour orchestrateur. Retourne
  \texttt{\{"status": "ok"\}} dès que l'application est lancée. \\
\hline
POST & \texttt{/recommend} &
  Calcule une recommandation. Corps requis~: \texttt{\{"code\_client":
  "<id>", "k": <1..3>\}}. Réponse~: \texttt{\{"mode":
  "model"|"cold\_start", "recommendations": [...]\}}. \\
\hline
\end{tabularx}
\caption{Endpoints HTTP exposés par l'API.}
\end{table}

\subsection{Modèle Pydantic de la Requête}

La validation entrée est assurée par un modèle Pydantic strict~:

\begin{lstlisting}[style=py, caption={Modèle Pydantic pour le corps de requête}]
from pydantic import BaseModel, Field

class RecommendRequest(BaseModel):
    code_client: str
    k: int = Field(default=1, ge=1, le=3)
\end{lstlisting}

L'attribut \texttt{Field(ge=1, le=3)} contraint $k$ à
$\{1, 2, 3\}$~; toute valeur hors borne déclenche une réponse HTTP~422
sans atteindre la logique métier.

\subsection{Exemples de Réponses}

\textbf{Cas modèle (utilisateur connu, $\geq 3$ réservations)~:}
\begin{lstlisting}[style=jn]
{
  "mode": "model",
  "recommendations": ["4512", "3801", "1122"]
}
\end{lstlisting}

\textbf{Cas cold start (utilisateur inconnu ou peu actif)~:}
\begin{lstlisting}[style=jn]
{
  "mode": "cold_start",
  "recommendations": []
}
\end{lstlisting}

\subsection{Démarrage et Vérification}

\begin{lstlisting}[style=sh]
.venv\Scripts\python.exe -m uvicorn apps.api.main:app --reload
# -> http://127.0.0.1:8000
# -> http://127.0.0.1:8000/docs   (Swagger UI auto-genere)
\end{lstlisting}

La documentation Swagger interactive est accessible à
\texttt{/docs}~; elle permet de tester les endpoints depuis le
navigateur sans recourir à un client externe.

\newpage
% ================================================================
\section{Tests Unitaires et Qualité du Code}
% ================================================================

\subsection{Stratégie de Test}

Le projet adopte une politique de tests systématiques sur les
composants à risque~: génération de candidats, ingénierie des
variables, calcul des métriques, logique du \texttt{Recommender},
fonctions d'entraînement. Les tests sont écrits avec \texttt{pytest}
et exécutés depuis l'environnement virtuel.

\subsection{Inventaire des Tests}

\begin{table}[h]
\centering
\small
\begin{tabularx}{\textwidth}{|l|c|X|}
\hline
\textbf{Fichier} & \textbf{Tests} & \textbf{Couverture} \\
\hline
\texttt{test\_candidates.py} & 11 &
  \texttt{generate\_candidates()}~: cas vide, cold start,
  tri par fréquence, tie-break par récence, fenêtre de lookback,
  dédoublonnage. \\
\hline
\texttt{test\_features.py}    & 12 &
  \texttt{build\_training\_rows()} (présence des 26 colonnes,
  types, encodage cyclique, \texttt{user\_top\_liaison\_share}
  $\in [0,1]$) et \texttt{compute\_inference\_row()} (1 ligne,
  \texttt{prev\_liaison} = dernier vu, \texttt{user\_trip\_index}
  = taille historique, \texttt{days\_since\_prev} $\geq 0$,
  features temporelles calées sur \texttt{asof}, levée
  d'exception sur historique vide). \\
\hline
\texttt{test\_metrics.py}     & 4  &
  \texttt{hit\_rate\_at\_k}, \texttt{mrr\_at\_k}~: cas trivial
  ($k=1$), cas multiclasse, cas dégénéré (proba uniforme). \\
\hline
\texttt{test\_recommender.py} & 4  &
  \texttt{recommend()}~: client connu actif, client cold start
  ($< 3$ réservations), client inconnu, contrainte sur $k$. \\
\hline
\texttt{test\_training.py}    & 2  &
  \texttt{temporal\_split} (proportion correcte, pas de
  chevauchement), \texttt{train\_xgb\_multiclass} (drop effectif
  de \texttt{CodeClient}). \\
\hline
\texttt{test\_cleaning.py}    & 5  &
  \texttt{make\_clean\_dataset()}~: propagation des annulations,
  filtrage cold start, jointure liaison, champs essentiels manquants,
  encodage cyclique sur cercle unité. \\
\hline
\texttt{test\_api.py}         & 5  &
  Endpoints HTTP via FastAPI \texttt{TestClient}~: \texttt{/health},
  \texttt{/recommend} mode modèle, mode cold start, validation
  Pydantic des bornes de $k$. \\
\hline
\textbf{Total}                & \textbf{43} & \\
\hline
\end{tabularx}
\caption{Inventaire de la suite de tests unitaires (état~:
43/43 verts).}
\end{table}

\subsection{Configuration \texttt{pytest}}

Le fichier \texttt{pyproject.toml} configure \texttt{pytest} avec
\texttt{pythonpath~=~["src"]}, ce qui rend les modules
\texttt{rec\_oncf} importables depuis les tests sans installation
préalable du package.

\textbf{Commande d'exécution~:}
\begin{lstlisting}[style=sh]
.venv\Scripts\python.exe -m pytest tests/ -v
\end{lstlisting}

\subsection{Conventions de Style et Modularité}

Le code source applique plusieurs disciplines~:
\begin{itemize}
\item \textbf{\texttt{from \_\_future\_\_ import annotations}}~: au
    sommet de chaque module pour profiter du typage différé Python~3.12.
\item \textbf{Type hints stricts}~: signatures de fonctions, classes
    \texttt{dataclass}, retours typés.
\item \textbf{Séparation données vs logique~:} le module
    \texttt{cleaning.py} ne dépend pas du module \texttt{training.py},
    et inversement.
\item \textbf{Pas de variables globales mutables}~: les artefacts sont
    passés explicitement entre fonctions.
\item \textbf{Erreurs explicites}~: \texttt{FileNotFoundError} avec
    message contextuel (\textit{« run scripts/02 first »}) plutôt que
    d'échouer silencieusement.
\end{itemize}

\newpage
% ================================================================
\section{Bilan et Perspectives}
% ================================================================

\subsection{Bilan Global du Projet}

À l'issue du projet, l'ensemble des objectifs initiaux est atteint~:
\begin{itemize}
\item \textbf{[OK]} Architecture deux étapes opérationnelle.
\item \textbf{[OK]} Métriques offline largement au-dessus des seuils
    (HR@1 = 76{,}3\,\% vs cible 30\,\%, après Sprint~2).
\item \textbf{[OK]} Conformité Loi 09-08 / CNDP respectée
    (\texttt{CodeClient} jamais utilisé comme feature).
\item \textbf{[OK]} API REST FastAPI prête à être consommée.
\item \textbf{[OK]} Reproductibilité complète garantie par les trois scripts
    numérotés et le fichier \texttt{pyproject.toml}.
\end{itemize}

L'API a été \textbf{validée en conditions réelles} en fin de phase~1~:
démarrage d'\texttt{uvicorn}, test de \texttt{/health}
($\rightarrow$~\texttt{\{"status": "ok"\}}), puis trois requêtes
\texttt{POST /recommend} couvrant le cas client très actif (679~voyages
historiques, top-3 retourné en mode \texttt{model}), client moyennement
actif (mode \texttt{model}), et client inconnu (mode
\texttt{cold\_start}). La validation Pydantic a été vérifiée sur
$k = 99$ (réponse HTTP~422 conforme).

Une mesure de latence a également été conduite (50 requêtes après
warmup) sur un poste portable~: \textbf{p50~=~852~ms}, p95~=~925~ms,
p99~=~975~ms. Cette latence \textbf{dépasse la cible SNF-02
($< 100$~ms)} : l'inférence multiclasse XGBoost sur 1\,011 classes
domine le temps de réponse. Une optimisation est planifiée en
phase~3 (modèle binaire pairwise scoré uniquement sur les 10
candidats au lieu d'une softmax sur 1\,011 classes).

\subsection{Travaux Sprint~2~: Industrialisation et Renforcement}

Une fois les jalons fonctionnels de la phase~1 atteints, un Sprint~2
a été conduit pour renforcer la qualité d'industrialisation du livrable.
Quatre chantiers ont été menés~:

\begin{enumerate}
\item \textbf{Métadonnées d'artefact.} Le sauveur de modèle dans
    \texttt{training.py} produit désormais un fichier
    \texttt{models/xgb\_ranker.meta.json} placé à côté du modèle, qui
    embarque~: date d'entraînement (UTC), versions exactes des
    bibliothèques (\texttt{xgboost~3.2.0}, \texttt{scikit-learn~1.8.0},
    \texttt{pandas~3.0.2}, \texttt{numpy~2.4.4}), volumétrie
    train/test, hyperparamètres, métriques offline, et un
    \texttt{dataset\_fingerprint} (SHA-256 tronqué de la forme et
    d'un échantillon du jeu de features). Ces métadonnées sont
    indispensables pour l'audit et la traçabilité en production~:
    elles permettent de reconstituer en quelle date un modèle a été
    formé, sur quelle version de données, avec quelles dépendances.

\item \textbf{Recalcul des features à la volée.} Les features
    fournies au modèle au moment d'une requête sont désormais calculées
    \textit{en ligne} à partir de l'historique frais du client, via
    \texttt{compute\_inference\_row()}. Le \texttt{Recommender}
    n'embarque plus le \texttt{feature\_lookup} statique pré-construit
    à l'entraînement, qui présentait le défaut critique de servir des
    features périmées dès qu'un client effectuait une nouvelle
    réservation. Ce changement débloque l'usage en production réelle
    --- sans lui, l'API ne pouvait fonctionner correctement qu'au
    moment de l'entraînement.

\item \textbf{Variable de fidélité \texttt{user\_top\_liaison\_share}.}
    Une nouvelle feature numérique a été ajoutée à la table
    (passant de 25 à 26 colonnes), capturant la concentration des
    voyages d'un client sur sa ligne dominante. Compatible avec
    l'architecture multiclasse, calculée sans fuite temporelle
    (cf.~section~6).

\item \textbf{Industrialisation logicielle.} Trois apports
    complémentaires~:
    (i)~un \texttt{.gitignore} excluant les artefacts générés
    (\texttt{models/}, \texttt{data/processed/}, \texttt{reports/},
    \texttt{.venv/})~;
    (ii)~un \texttt{README.md} synthétique en porte d'entrée du dépôt~;
    (iii)~un \textbf{workflow GitHub Actions}
    (\texttt{.github/workflows/tests.yml}) qui exécute
    \texttt{pytest} et \texttt{ruff} à chaque \textit{push} ou
    \textit{pull request}, garantissant que la suite de tests
    (43 tests) reste verte sans intervention manuelle.
\end{enumerate}

\subsection{État d'Avancement Détaillé}

\begin{table}[h]
\centering
\small
\begin{tabularx}{\textwidth}{|l|c|X|}
\hline
\textbf{Composant} & \textbf{État} & \textbf{Notes} \\
\hline
Nettoyage des données          & Terminé & 491\,680 lignes propres en sortie \\
\hline
Feature engineering            & Terminé & 26 colonnes, 491\,680 lignes \\
\hline
Entraînement XGBoost           & Terminé & Métriques au-dessus des seuils \\
\hline
Composant \texttt{Recommender} & Terminé & Encapsule toute la logique \\
\hline
Suite de tests (43 tests)      & Terminé & 43/43 verts \\
\hline
API FastAPI (code)             & Terminé & Modèle disponible, prêt à démarrer \\
\hline
API en conditions réelles      & Terminé & 5 cas testés en live (health + 3 recommend + validation k) \\
\hline
Latence p50/p95 mesurée        & Mesurée & 852~ms / 925~ms (au-dessus de SNF-02, optim phase~3) \\
\hline
Baselines de comparaison       & Terminé & most\_frequent, prev\_liaison, global\_top calculés \\
\hline
Métadonnées d'artefact         & Terminé & \texttt{models/xgb\_ranker.meta.json} (versions, hyperparams, fingerprint) \\
\hline
Recalcul features en ligne     & Terminé & \texttt{compute\_inference\_row()} remplace le \texttt{feature\_lookup} stale \\
\hline
Feature \texttt{user\_top\_liaison\_share} & Terminé & 26 colonnes au total \\
\hline
.gitignore + README + CI       & Terminé & GitHub Actions exécute \texttt{pytest} + \texttt{ruff} sur push/PR \\
\hline
Intégration planning ONCF      & Déféré & Phase~3~: enrichissement temps réel \\
\hline
Réentraînement automatisé      & Déféré & Phase~2~: pipeline CI/CD \\
\hline
A/B test en production         & Déféré & Phase~2~: validation business \\
\hline
\end{tabularx}
\caption{État d'avancement par composant.}
\end{table}

\subsection{Perspectives d'Amélioration (Phase~3)}

Plusieurs axes d'évolution structurent la feuille de route post-PFA.

\textbf{Axe~1~: Refonte en \textit{pairwise ranker} (priorité
maximale).} La latence actuelle de l'API ($\approx 850$~ms p50) est
dominée par le coût d'inférence du multiclasse softmax sur 1\,011
classes. La refonte en \textit{pairwise ranker} (\texttt{XGBRanker}
avec \texttt{rank:pairwise} ou modèle binaire scoré par candidat)
réduirait ce coût d'environ deux ordres de grandeur~: scorer 10
candidats au lieu de 1\,011 classes signifie $\approx 100\times$ moins
d'opérations par requête. Cette refonte débloque également
l'utilisation de \textit{features par candidat} (popularité globale
par liaison, score de similarité O/D), incompatibles avec la
formulation multiclasse actuelle.

\textbf{Axe~2~: Tuning hyperparamétrique systématique.} Une recherche
\texttt{Optuna} sur l'espace des hyperparamètres
(\texttt{n\_estimators}, \texttt{max\_depth}, \texttt{learning\_rate},
\texttt{reg\_lambda}, \texttt{subsample}) pourrait révéler des
configurations supérieures. La contrainte CPU (43~min/run) limite
cependant le nombre d'essais possibles~; un GPU plus capable
(8~Go+ VRAM) ouvrirait cette possibilité.

\textbf{Axe~3~: Réentraînement automatisé avec garde-fous.} Un
pipeline CI/CD hebdomadaire qui (i)~régénère
\texttt{features.parquet}, (ii)~réentraîne le modèle, (iii)~le promeut
en production uniquement si \texttt{HR@1} ne baisse pas de plus de
5~points par rapport au champion, (iv)~archive l'ancien modèle pour
\textit{rollback} en cas d'incident. La métadonnée d'artefact
introduite en Sprint~2 facilite directement ce mécanisme de
gouvernance.

\textbf{Axe~4~: Logging structuré et observabilité.} Ajouter à l'API
un logger \texttt{structlog} émettant pour chaque requête~: latence,
mode (\texttt{model} / \texttt{cold\_start}), version du modèle utilisé.
Le \texttt{code\_client} ne doit jamais apparaître en clair (Loi~09-08).
Connexion à un dashboard temps réel (Grafana, Datadog) pour piloter
le \texttt{cold\_start\_rate}, la latence p95 et la dérive de
distribution des prédictions.

\textbf{Axe~5~: Intégration avec le planning ONCF.} Enrichir les
recommandations par les disponibilités effectives des trains permettrait
de présenter directement à l'utilisateur les horaires concrets, plutôt
que de simples identifiants de liaison. Ce chantier requiert une API
fournie par la Direction Digitale ONCF.

\textbf{Axe~6~: A/B testing en production.} La validation finale du
gain produit doit être faite en ligne~: un échantillon d'utilisateurs
reçoit les recommandations, un autre conserve l'expérience classique~;
le différentiel de CTR sur l'écran d'accueil mesure le gain réel.
Endpoint suggéré~: \texttt{/recommend?variant=A|B} avec routage
déterministe par hash du \texttt{code\_client}.

\textbf{Axe~7~: Embeddings de liaisons (long terme).} À mesure que le
volume de données augmente, l'apprentissage d'embeddings de liaisons
(à la manière d'Airbnb, via Word2Vec sur séquences de trajets) deviendrait
viable et pourrait capturer des similarités sémantiques entre trajets
(gares jumelles, axes parallèles, saisonnalité partagée).

\subsection{Apport Personnel}

Au-delà des livrables techniques, ce projet a constitué une expérience
formatrice sur plusieurs plans~:
\begin{itemize}
\item \textbf{Ingénierie data réaliste~:} confrontation à des données
    brutes réelles, avec leur lot de bizarreries (annulations
    masquées, valeurs parasites, colonnes constantes).
\item \textbf{Discipline méthodologique~:} importance du split temporel
    correct, du cold start strict, de la séparation logique
    métier/transport.
\item \textbf{Sensibilité réglementaire~:} application concrète de
    principes de protection des données (Loi 09-08), au-delà de la
    seule théorie.
\item \textbf{Dépannage de problèmes opaques~:} l'épisode de l'OOM
    silencieux sur GPU a illustré combien le diagnostic peut être
    difficile lorsqu'il n'y a aucun message d'erreur exploitable.
\item \textbf{Industrialisation légère~:} encapsulation propre,
    tests unitaires, documentation, scripts numérotés
    reproductibles --- des bonnes pratiques transférables à tout futur
    contexte professionnel.
\end{itemize}

\newpage
% ================================================================
\section{Conclusion Générale}
% ================================================================

Ce Projet de Fin d'Année a permis la réalisation, de bout en bout,
d'un \textbf{système de recommandation proactif} pour l'opérateur
ferroviaire ONCF, depuis la donnée brute jusqu'à l'API REST prête à
être consommée par l'application mobile.

Sur le plan \textbf{conceptuel}, le projet a confronté des choix
d'architecture (\textit{deux étapes Candidate~+~Ranking}, modèle
\textit{multiclasse} versus \textit{learning to rank}, \textit{cold
start} strict vs heuristique de secours) à un cas d'usage
industriel concret. Le benchmark préalable des pratiques d'Uber, de
Trainline, d'Airbnb et de Google Maps a guidé ces choix vers des
solutions éprouvées.

Sur le plan \textbf{technique}, le pipeline complet --- nettoyage de
946\,155 lignes brutes, ingénierie de 26 features comportementales et
temporelles, entraînement XGBoost multiclasse sur 1\,011 classes,
encapsulation dans un composant \texttt{Recommender} consommable par
une API FastAPI --- est opérationnel et testé. Les métriques offline
atteintes après Sprint~2 (HR@1 = 76{,}3\,\%, HR@3 = 90{,}6\,\%, MRR@3 = 82{,}8\,\%)
dépassent significativement les seuils cibles initiaux et confirment
la pertinence des choix de conception. La progression monotone des
métriques par segment d'historique (de 73{,}9\,\% pour les voyageurs
ayant moins de trois voyages dans le train à \textbf{82{,}7\,\%} pour les
fidèles avec plus de vingt voyages) valide l'hypothèse que le modèle
apprend des patterns spécifiquement individuels.

Sur le plan \textbf{réglementaire}, la conformité à la Loi~09-08 est
garantie par construction~: le \texttt{CodeClient} reste une clé de
lookup et n'est jamais utilisé comme variable d'apprentissage. Les
incidents historiques (fuite identifiée et corrigée pendant le
développement) ont fait l'objet d'un test unitaire défensif pour éviter
toute régression future.

Sur le plan \textbf{personnel}, ce projet a constitué une immersion
significative dans les pratiques d'ingénierie data en environnement
réel~: gestion de données imparfaites, dépannage de pannes opaques
(saturation VRAM silencieuse), discipline de validation temporelle,
séparation des préoccupations, automatisation par tests.

Les \textbf{perspectives à court terme} (validation API, ajout de
features de popularité globale, tuning hyperparamétrique) et à
\textbf{moyen terme} (réentraînement automatisé, A/B testing en ligne,
intégration du planning ONCF temps réel) tracent une feuille de route
claire pour transformer ce livrable de PFA en système de production
durable, capable d'apporter une valeur tangible aux millions de
voyageurs qui utilisent quotidiennement les services de l'ONCF.

\bigskip

Au final, ce projet illustre comment des techniques de
\textit{machine learning} bien maîtrisées --- ni sur-dimensionnées ni
sous-équipées --- peuvent traiter avec efficacité un problème
métier concret, en respectant les contraintes opérationnelles et
réglementaires d'une grande organisation publique. Le fait que la
solution puisse être reproduite à partir de trois scripts
numérotés et d'une commande \texttt{uvicorn}, sur un ordinateur
portable et en quelques heures, est sans doute le meilleur indicateur
de sa robustesse de conception.

\newpage
% ================================================================
\section*{Bibliographie / Webographie}
\addcontentsline{toc}{section}{Bibliographie / Webographie}
% ================================================================

\begin{enumerate}
\item Chen,~T., Guestrin,~C. (2016). \textit{XGBoost~: A Scalable
    Tree Boosting System}. Proceedings of the 22nd ACM SIGKDD
    International Conference on Knowledge Discovery and Data Mining
    (KDD).
\item Grbovic,~M., Cheng,~H. (2018). \textit{Real-time Personalization
    using Embeddings for Search Ranking at Airbnb}. KDD~2018.
\item Documentation officielle XGBoost~3.x~:
    \url{https://xgboost.readthedocs.io/}.
\item Documentation officielle scikit-learn~:
    \url{https://scikit-learn.org/stable/}.
\item Documentation officielle FastAPI~:
    \url{https://fastapi.tiangolo.com/}.
\item Pydantic~v2 documentation~: \url{https://docs.pydantic.dev/}.
\item Site officiel de l'ONCF~: \url{https://www.oncf.ma/}.
\item Loi 09-08 relative à la protection des personnes physiques à
    l'égard du traitement des données à caractère personnel
    (Royaume du Maroc, 2009).
\item Site de la Commission Nationale de contrôle de la protection des
    Données à caractère Personnel (CNDP)~:
    \url{https://www.cndp.ma/}.
\item Uber Engineering Blog~: \textit{Personalized Recommendations at
    Uber Eats}, articles 2018--2022.
\item Karpathy,~A. (2019). \textit{A Recipe for Training Neural
    Networks}, blog post (méthodologie générale appliquée par analogie
    à XGBoost).
\end{enumerate}

\newpage
% ================================================================
\appendix
\section{Annexe~A~: Rapport de Nettoyage Complet}
% ================================================================

\begin{lstlisting}[style=jn]
{
  "rows_before": 946155,
  "rows_after": 491680,
  "cold_start_removed_rows": 403615,
  "cold_start_removed_clients": 304595,
  "distinct_clients": 69449,
  "distinct_liaisons": 1067,
  "join_overlap_rate": 0.9992,
  "min_trips_for_training": 3
}
\end{lstlisting}

\section{Annexe~B~: Métriques Offline Détaillées (Sprint~2)}
\begin{lstlisting}[style=jn]
{
  "rows_train": 393344,
  "rows_test": 98261,
  "rows_test_unseen_dropped": 75,
  "classes": 1011,
  "n_estimators": 200,
  "max_depth": 6,
  "device": "cpu",
  "hit_rate@1": 0.7627746511840914,
  "hit_rate@3": 0.9055474705122073,
  "mrr@3": 0.8276885030683587,
  "thresholds_met": {
    "hit_rate@1 > 0.30": true,
    "hit_rate@3 > 0.50": true,
    "mrr@3 > 0.35": true
  },
  "metrics_by_history_segment": {
    "0_2_unseen_in_train": {
      "n": 44397, "hit_rate@1": 0.7393, "hit_rate@3": 0.8930, "mrr@3": 0.8091
    },
    "3_5_new_active": {
      "n": 16780, "hit_rate@1": 0.7517, "hit_rate@3": 0.9038, "mrr@3": 0.8210
    },
    "6_20_regular": {
      "n": 23737, "hit_rate@1": 0.7786, "hit_rate@3": 0.9159, "mrr@3": 0.8411
    },
    "21_plus_loyal": {
      "n": 13347, "hit_rate@1": 0.8268, "hit_rate@3": 0.9311, "mrr@3": 0.8741
    }
  }
}
\end{lstlisting}

\section{Annexe~B'~: Métadonnées du Modèle (\texttt{xgb\_ranker.meta.json})}
\begin{lstlisting}[style=jn]
{
  "trained_at_utc": "2026-05-04T10:58:06+00:00",
  "platform": "Windows-11-10.0.26200-SP0",
  "python": "3.12.10",
  "package_versions": {
    "xgboost": "3.2.0",
    "scikit-learn": "1.8.0",
    "pandas": "3.0.2",
    "numpy": "2.4.4"
  },
  "train_rows": 393344, "test_rows": 98261, "n_classes": 1011,
  "hyperparameters": {
    "objective": "multi:softprob", "n_estimators": 200, "max_depth": 6,
    "learning_rate": 0.08, "subsample": 0.9, "colsample_bytree": 0.8,
    "reg_lambda": 1.0, "tree_method": "hist", "device": "cpu",
    "random_state": 42
  },
  "metrics": {
    "hit_rate@1": 0.7628, "hit_rate@3": 0.9055, "mrr@3": 0.8277
  },
  "dataset_fingerprint": "4d0dfd12e0b60341"
}
\end{lstlisting}

\section{Annexe~C~: Code Source du Script d'Entraînement}

\begin{lstlisting}[style=py, caption={\texttt{scripts/03\_train\_ranker.py}}]
from __future__ import annotations
import json, sys
from pathlib import Path

PROJECT_ROOT = Path(__file__).resolve().parents[1]
sys.path.insert(0, str(PROJECT_ROOT / "src"))

from rec_oncf.config import default_paths
from rec_oncf.io import read_parquet
from rec_oncf.metrics import hit_rate_at_k, mrr_at_k
from rec_oncf.training import (
    load_artifacts, predict_proba, save_artifacts,
    temporal_split, train_xgb_multiclass,
)

def main():
    paths = default_paths()
    df = read_parquet(paths.features_parquet).drop(columns=["CodeClient"])

    label_col = "LiaisonId"
    time_col  = "DateHeureDepartVoyageSegment"

    train_df, test_df = temporal_split(df, time_col=time_col, train_frac=0.8)
    artifacts = train_xgb_multiclass(train_df, label_col=label_col, time_col=time_col)
    save_artifacts(artifacts, model_path=paths.xgb_model_path,
                   label_encoder_path=paths.label_encoder_path)

    artifacts = load_artifacts(model_path=paths.xgb_model_path,
                               label_encoder_path=paths.label_encoder_path)

    known = set(artifacts.label_encoder.classes_)
    test_df = test_df[test_df[label_col].astype(str).isin(known)]

    proba = predict_proba(artifacts, test_df, label_col=label_col)
    y_true = artifacts.label_encoder.transform(test_df[label_col].astype(str).to_numpy())

    hr1  = hit_rate_at_k(y_true, proba, k=1)
    hr3  = hit_rate_at_k(y_true, proba, k=3)
    mrr3 = mrr_at_k(y_true, proba, k=3)

    report = {
        "rows_train": int(len(train_df)),
        "rows_test":  int(len(test_df)),
        "classes":    int(len(artifacts.label_encoder.classes_)),
        "hit_rate@1": float(hr1),
        "hit_rate@3": float(hr3),
        "mrr@3":      float(mrr3),
    }

    out = paths.project_root / "reports" / "offline_metrics.json"
    out.parent.mkdir(parents=True, exist_ok=True)
    out.write_text(json.dumps(report, indent=2), encoding="utf-8")
    print(json.dumps(report, indent=2))

if __name__ == "__main__":
    main()
\end{lstlisting}

\section{Annexe~D~: Commandes de Reproduction Complète}

\begin{lstlisting}[style=sh]
# Prerequis : oncf_data.csv et Liaison.csv places sur le Bureau
python -m venv .venv
.venv\Scripts\pip install -e .

# 1. Nettoyage -> data/processed/oncf_clean.parquet
.venv\Scripts\python scripts\01_make_dataset.py

# 2. Features -> data/processed/features.parquet (26 colonnes)
.venv\Scripts\python scripts\02_build_features.py

# 3. Entrainement XGBoost (~43 min CPU)
#    -> models/xgb_ranker.json
#    -> models/label_encoder.joblib
#    -> reports/offline_metrics.json
.venv\Scripts\python scripts\03_train_ranker.py

# 4. Tests unitaires (26 tests)
.venv\Scripts\python -m pytest tests/ -v

# 5. Demarrage de l'API REST
.venv\Scripts\python -m uvicorn apps.api.main:app --reload
# -> http://127.0.0.1:8000
# -> http://127.0.0.1:8000/docs   (Swagger UI)
\end{lstlisting}

\end{document}
